% Relatório técnico --- Detecção de anomalias em telemetria com redes de Hopfield
% Compilador: pdfLaTeX (padrão do Overleaf). Documento principal: main.tex
\documentclass[a4paper,11pt]{article}
% ===================== Codificação, idioma e tipografia =====================
\usepackage[utf8]{inputenc}
\usepackage[T1]{fontenc}
\usepackage{lmodern}
\usepackage[brazilian]{babel}
\usepackage{csquotes}
\usepackage{microtype}
\usepackage[a4paper,margin=2.5cm]{geometry}
\usepackage{setspace}
\setstretch{1.1}
\setlength{\emergencystretch}{2em}

% ============================== Matemática ==================================
\usepackage{amsmath,amssymb,amsthm,mathtools}
\numberwithin{equation}{section}

% ====================== Tabelas, figuras, listas e caixas ===================
\usepackage{graphicx}
\usepackage{booktabs,array,tabularx,longtable,multirow}
\usepackage{xcolor}
\usepackage[font=small,labelfont=bf,labelsep=period]{caption}
\usepackage{float}
\usepackage{enumitem}
\setlist{itemsep=2pt,topsep=4pt}
\usepackage{tikz}
\usetikzlibrary{arrows.meta,positioning,calc,fit,shapes.geometric}
\usepackage[skins,breakable]{tcolorbox}

\newcolumntype{Y}{>{\centering\arraybackslash}X}
\newcolumntype{L}{>{\raggedright\arraybackslash}X}
\newcolumntype{R}{>{\raggedleft\arraybackslash}X}
\newcolumntype{P}[1]{>{\raggedright\arraybackslash}p{#1}}
\newtcolorbox{notabox}[1]{colback=yellow!8,colframe=orange!60!black,
  fonttitle=\bfseries,title={#1},boxrule=0.6pt,arc=2pt,breakable}

% =========================== Algoritmos e listagens =========================
\usepackage{algorithm}
\usepackage{algpseudocode}
\floatname{algorithm}{Algoritmo}
\algrenewcommand\algorithmicrequire{\textbf{Entrada:}}
\algrenewcommand\algorithmicensure{\textbf{Saída:}}
\algrenewcommand\algorithmicend{\textbf{fim}}
\algrenewcommand\algorithmicwhile{\textbf{enquanto}}
\algrenewcommand\algorithmicdo{\textbf{faça}}
\algrenewcommand\algorithmicfor{\textbf{para}}
\algrenewcommand\algorithmicforall{\textbf{para todo}}
\algrenewcommand\algorithmicif{\textbf{se}}
\algrenewcommand\algorithmicthen{\textbf{então}}
\algrenewcommand\algorithmicelse{\textbf{senão}}
\algrenewcommand\algorithmicreturn{\textbf{retorne}}
\algrenewcommand\algorithmicrepeat{\textbf{repita}}
\algrenewcommand\algorithmicuntil{\textbf{até que}}

\usepackage{listings}
\lstset{basicstyle=\ttfamily\footnotesize,breaklines=true,frame=single,
  columns=fullflexible,keepspaces=true,xleftmargin=2pt,xrightmargin=2pt}

% ============================ Teoremas e definições =========================
\theoremstyle{plain}
\newtheorem{proposicao}{Proposição}[section]
\theoremstyle{definition}
\newtheorem{definicao}[proposicao]{Definição}
\theoremstyle{remark}
\newtheorem{observacao}[proposicao]{Observação}
\addto\captionsbrazilian{\renewcommand{\proofname}{Demonstração}}

% ================================= Macros ===================================
\newcommand{\R}{\mathbb{R}}
\DeclareMathOperator{\sgn}{sgn}
\DeclareMathOperator*{\argmin}{arg\,min}
\DeclareMathOperator*{\argmax}{arg\,max}
\DeclareMathOperator{\lse}{lse}
\DeclareMathOperator{\softmax}{softmax}
\newcommand{\hamm}{d_{\mathrm{H}}}
\newcommand{\sstar}{s^{\ast}}
% Marcadores de conteúdo pendente (texto e células de tabela)
\newcommand{\apreencher}[1][]{\textcolor{red!70!black}{\textbf{[a preencher\ifx&#1&\else: #1\fi]}}}
\newcommand{\lacuna}{\textcolor{red!70!black}{$\cdots$}}

% ======================= Hiperlinks (carregar por último) ===================
\usepackage[colorlinks=true,linkcolor=blue!50!black,citecolor=blue!50!black,
  urlcolor=blue!50!black,unicode=true,
  pdftitle={Detecção de Anomalias em Telemetria de Sistemas Utilizando Redes de Hopfield},
  pdfauthor={Equipe}]{hyperref}
\usepackage{xurl}


\title{\textbf{Detecção de Anomalias em Telemetria de Sistemas\\ Utilizando Redes de Hopfield}\\[0.8em]
  \large Trabalho de Implementação de Redes Neurais Artificiais\\
  \large Tema: Redes baseadas em Memória Associativa (Redes Neurais Recorrentes)}
\author{\apreencher[nome do integrante 1] \and \apreencher[nome do integrante 2] \and
  \apreencher[nome do integrante 3]\\[0.6em]
  \small \apreencher[instituição, curso e disciplina]}
\date{\today}

\begin{document}
\maketitle

\begin{abstract}
Este relatório apresenta o projeto de um detector de anomalias em telemetria de sistemas baseado em
redes de Hopfield, o modelo clássico de memória associativa recorrente. O problema é tratado como
classificação binária (normal $\times$ anômalo) sobre métricas de utilização e atividade de máquinas
servidoras (carga de CPU, uso de memória e de rede, entre outras). Usamos o \emph{Server Machine Dataset}
(SMD): 28 máquinas, 38 métricas amostradas por minuto e 708.420 pontos de teste, dos quais 4,16\,\% são
anômalos. A abordagem armazena na matriz de pesos protótipos do funcionamento normal (regras de Hebb e de
Storkey), recupera cada observação por dinâmica assíncrona e usa distância de Hamming e energia como escores
de anomalia; a rede de Hopfield moderna (Ramsauer et al.) entra como variante. O texto cobre a teoria
(energia, convergência, capacidade, estados espúrios), a análise dos dados, o protocolo experimental
(divisões cronológicas, métricas, linhas de base clássicas e o protocolo com LLM da Etapa~II) e uma
comparação crítica a priori. Os resultados experimentais ainda não foram obtidos: as tabelas
correspondentes estão marcadas como lacunas a preencher.

\medskip
\noindent\textbf{Palavras-chave:} redes de Hopfield; memória associativa; detecção de anomalias;
telemetria de sistemas; classificação binária; \emph{Server Machine Dataset}.
\end{abstract}

\thispagestyle{empty}
\newpage
\setcounter{tocdepth}{2}
\tableofcontents
\newpage

\section{Introdução}\label{sec:intro}

\subsection{Contexto e motivação}
Sistemas operacionais e servidores expõem continuamente contadores de utilização e de atividade em baixo
nível --- carga de CPU, uso de memória, tráfego de rede, entre outros. Essa \emph{telemetria} é a base do
monitoramento: quando o comportamento de uma máquina se afasta do habitual (um vazamento de memória, um
processo descontrolado, uma falha de rede), o desvio aparece nas métricas antes de se tornar uma
interrupção visível. Detectar automaticamente esses desvios é o problema clássico de \emph{detecção de
anomalias}~\cite{chandola2009}.

A dificuldade prática é dupla. Primeiro, anomalias são raras e heterogêneas: nos dados deste trabalho
apenas 4,16\,\% dos pontos de teste são anômalos (Seção~\ref{sec:dados}), de modo que é impossível
enumerar de antemão todas as formas de falha. Segundo, o rótulo ``anômalo'' costuma vir de relatórios de
incidentes, ao passo que o comportamento \emph{normal} está disponível em abundância e sem rótulo. Isso
favorece abordagens que \emph{aprendem o que é normal} e tratam como anomalia aquilo que não se encaixa.

As redes de Hopfield~\cite{hopfield1982} são um ajuste natural a essa ideia. Trata-se de uma rede neural
recorrente totalmente conectada, cujo estado evolui, sob pesos simétricos, até um mínimo local de uma função
de energia. Os padrões gravados na matriz de pesos tornam-se atratores, e apresentar à rede uma entrada
parcial ou ruidosa a faz convergir para o padrão armazenado mais próximo~\cite{hopfieldtank1985}. Se os
padrões armazenados descrevem o funcionamento normal, esperamos que uma observação normal seja
``corrigida'' para um padrão familiar com pouco esforço, enquanto uma observação anômala permaneça distante
de qualquer atrator ou seja atraída para um estado que não corresponde a nenhuma memória. Essa diferença na
recuperação pode ser convertida em um escore de anomalia. A ideia é uma hipótese de trabalho, não um
resultado: avaliá-la é o objetivo do projeto.

\subsection{Objetivos}
\paragraph{Objetivo geral.} Avaliar a capacidade de uma rede de Hopfield de armazenar padrões de
funcionamento normal e reconhecer comportamentos que se desviam desses padrões, tratando a detecção de
anomalias em telemetria de sistemas como um problema de classificação binária (normal $\times$ anômalo).

\paragraph{Objetivos específicos.}
\begin{enumerate}
  \item Revisar o modelo de Hopfield em suas formas discreta, contínua (Hopfield--Tank) e moderna
        (Ramsauer et al.), com foco em energia, convergência, capacidade e estados espúrios.
  \item Formalizar um detector \emph{one-class} baseado em recuperação associativa, com escores definidos a
        partir de distância de Hamming e de energia.
  \item Caracterizar o conjunto de dados (SMD): amostras, atributos, rótulos, desbalanceamento e deriva.
  \item Definir o protocolo experimental: preparação, codificação binária, divisões cronológicas, métricas
        e linhas de base clássicas (Etapa~I).
  \item Definir o protocolo da Etapa~II (LLM com \emph{zero-shot} e \emph{few-shot}) sobre os mesmos dados.
  \item Formular hipóteses verificáveis e critérios para a comparação crítica dos resultados.
\end{enumerate}

\subsection{Pergunta de pesquisa}
\emph{Uma memória associativa de Hopfield, treinada somente com telemetria predominantemente normal,
separa pontos normais de pontos anômalos? E como esse desempenho se compara ao de classificadores clássicos
e ao de um LLM usado por \emph{prompting}?}

\subsection{Escopo e premissas}
\begin{notabox}{Escopo desta versão do relatório}
\begin{itemize}
  \item \textbf{Projeto e fundamentação, sem experimentos.} Este texto não contém resultados de modelos.
        As tabelas da Seção~\ref{sec:resultados} têm células vazias (\lacuna) a preencher após os
        experimentos; nenhum valor foi inventado. Os números das Seções~\ref{sec:dados} e~\ref{sec:metodo}
        são estatísticas reais do SMD (contagens, divisões e exploração dos dados) e estão identificados
        como tais.
  \item \textbf{Hopfield não consta da lista do enunciado.} Para classificação, o enunciado cita \emph{Decision
        Trees}, \emph{Random Forest}, \emph{Naive Bayes}, SVM, MLP, XGBoost e KNN~\cite{enunciado}. Adotamos
        a rede de Hopfield como \emph{modelo do tema} e \emph{Random Forest} e SVM como linhas de base
        (mínimo de 2 modelos da Etapa~I). \textbf{Confirmar este arranjo com o professor.}
  \item \textbf{``Baixo nível'' no SMD.} São 38 contadores anonimizados de utilização de recursos, sem
        chamadas de sistema nem rastros de \emph{kernel}. Coletar telemetria própria é trabalho futuro
        (Seção~\ref{sec:conclusoes}).
  \item \textbf{Etapa~II é um protocolo.} A execução exige acesso a um modelo de linguagem; aqui
        documentamos tarefa, \emph{prompt}, amostragem e métricas (Apêndice~\ref{apx:prompt}).
  \item \textbf{Propostas, não resultados.} Divisões, hiperparâmetros e linhas de base são decisões de
        projeto a validar; as derivações próprias estão sinalizadas como tais.
\end{itemize}
\end{notabox}

\subsection{Organização do texto}
A Seção~\ref{sec:teoria} apresenta a fundamentação teórica. A Seção~\ref{sec:dados} define o problema e
descreve o conjunto de dados, e a Seção~\ref{sec:metodo} detalha a metodologia, incluindo as Etapas~I e~II.
A Seção~\ref{sec:resultados} reserva as tabelas de resultados, a Seção~\ref{sec:critica} traz a comparação
crítica (hipóteses e pontos positivos e negativos) e a Seção~\ref{sec:conclusoes} fecha com conclusões e
próximos passos. Os apêndices reúnem pseudocódigo, um exemplo numérico, o \emph{prompt}, o cronograma, o
mapeamento dos requisitos do enunciado, derivações e a tabela das 28 máquinas.

\section{Fundamentação teórica}\label{sec:teoria}

Usamos $N$ para o número de neurônios (dimensão do estado) e $K$ para o número de padrões armazenados.%
\footnote{Ramsauer et al.~\cite{ramsauer2021} usam $d$ para a dimensão e $N$ para o número de padrões; as
fórmulas deles são reescritas aqui na nossa notação.}
As fontes principais são a bibliografia indicada pelo professor --- Hopfield e Tank~\cite{hopfieldtank1985},
o verbete da Wikipédia~\cite{wikihopfield} e Ramsauer et al.~\cite{ramsauer2021} --- e, quando necessário, as
obras originais por elas citadas. Derivações próprias estão sinalizadas como tais.

\subsection{Memória associativa e redes recorrentes}
Uma memória endereçável por conteúdo (\emph{content-addressable memory}, CAM) recupera um item completo a
partir de uma parte dele ou de uma versão ruidosa, em vez de acessá-lo por endereço. Hopfield~\cite{hopfield1982}
mostrou que uma rede recorrente de elementos de limiar com conexões simétricas realiza essa função: a
recuperação é a evolução dinâmica de um estado inicial (a ``chave'') até um estado estável, interpretado como
a memória evocada. Em Hopfield e Tank~\cite[p.~144]{hopfieldtank1985}, cada recuperação ``resolve o problema''
de qual das palavras binárias armazenadas é a mais próxima da chave. É essa propriedade --- identificar o
padrão normal mais próximo --- que exploramos como detector de anomalias.

\subsection{A rede de Hopfield discreta}\label{sec:discreta}
\begin{definicao}[Rede de Hopfield discreta]
Uma rede de Hopfield discreta com $N$ neurônios é definida por uma matriz de pesos simétrica
$W\in\R^{N\times N}$ ($w_{ij}=w_{ji}$) com diagonal nula ($w_{ii}=0$) e por limiares $\theta\in\R^{N}$. O
estado é $s\in\{-1,+1\}^{N}$, o campo local do neurônio $i$ é $h_i(s)=\sum_{j}w_{ij}s_j$ e sua atualização é
\begin{equation}\label{eq:update}
s_i\leftarrow\begin{cases}+1,& h_i(s)\ge\theta_i,\\ -1,& \text{caso contrário.}\end{cases}
\end{equation}
\end{definicao}

Essa é a formulação da Wikipédia~\cite{wikihopfield}: unidades $\pm1$, sem autoconexões, pesos simétricos e
limiares em geral nulos e não treinados. A rede é um grafo completo não direcionado
(Figura~\ref{fig:grafo}). A atualização pode ser \emph{assíncrona} (um neurônio por vez, em ordem aleatória
ou fixa) ou \emph{síncrona} (todos de uma vez, o que exige um relógio central)~\cite{wikihopfield}. A cada
estado associa-se a energia
\begin{equation}\label{eq:energia}
E(s)=-\tfrac12\sum_{i,j}w_{ij}s_is_j+\sum_i\theta_is_i=-\tfrac12\,s^{\top}Ws+\theta^{\top}s .
\end{equation}

\begin{figure}[htbp]
\centering
\begin{tikzpicture}[neuron/.style={circle,draw,thick,minimum size=8mm,fill=blue!8,inner sep=0pt}]
  \foreach \i in {1,...,5}{
    \node[neuron] (n\i) at ({90+72*(\i-1)}:1.7) {$s_{\i}$};
  }
  \foreach \i in {1,...,4}{
    \foreach \j in {\the\numexpr\i+1\relax,...,5}{
      \draw[gray!80] (n\i) -- (n\j);
    }
  }
  \node[align=left,anchor=west,font=\small] at (3.0,0) {$N=5$ neurônios\\[2pt]
    $\tfrac{N(N-1)}{2}=10$ conexões\\[2pt] $w_{ij}=w_{ji},\ w_{ii}=0$};
\end{tikzpicture}
\caption{Rede de Hopfield com $N=5$ neurônios: grafo completo não direcionado, sem autoconexões.}
\label{fig:grafo}
\end{figure}

\begin{proposicao}[Convergência assíncrona]\label{prop:conv}
Se $W$ é simétrica com diagonal nula, cada atualização assíncrona não aumenta $E$ e o número de mudanças de
estado é finito. Assim, varreduras sucessivas (em ordem fixa ou aleatória) terminam em um ponto fixo, que é
um mínimo local de $E$ em relação a mudanças de um único neurônio.
\end{proposicao}
\begin{proof}
Seja $k$ o neurônio atualizado e $s'$ o estado resultante. Como $w_{kk}=0$ e $W$ é simétrica, apenas os
termos de $E$ que envolvem $s_k$ mudam, e
\[
E(s')-E(s)=-(s_k'-s_k)\bigl(h_k(s)-\theta_k\bigr).
\]
Se $s_k'\neq s_k$ e $h_k>\theta_k$, então $s_k'=+1$, $s_k=-1$ e a variação é $-2(h_k-\theta_k)<0$; se
$h_k<\theta_k$, então $s_k'=-1$, $s_k=+1$ e a variação é $2(h_k-\theta_k)<0$. O único caso sem queda estrita
é o empate $h_k=\theta_k$, em que a regra~\eqref{eq:update} só permite a transição $-1\to+1$. Como há $2^N$
estados, $E$ assume um número finito de valores e só cai estritamente um número finito de vezes. Entre duas
quedas estritas consecutivas ocorrem no máximo $N$ transições por empate, pois cada neurônio só pode ir de
$-1$ para $+1$ uma vez sem que haja queda estrita. Logo o número total de mudanças é finito.
\end{proof}

\begin{observacao}[Dinâmica síncrona]
Com atualização síncrona a monotonicidade de $E$ não vale. Para $W=\bigl(\begin{smallmatrix}0&-1\\-1&0\end{smallmatrix}\bigr)$
e $\theta=0$, o estado $(+1,+1)$ vai para $(-1,-1)$ e volta, num ciclo de período~2, ao passo que a
atualização assíncrona leva a $(-1,+1)$, que é ponto fixo. Bruck~\cite{bruck1990} provou um teorema de
convergência generalizado que cobre as dinâmicas assíncrona e síncrona~\cite{wikihopfield}. Neste trabalho
usamos a dinâmica assíncrona.
\end{observacao}

\subsection{Armazenamento: regras de Hebb e de Storkey}\label{sec:regras}
\paragraph{Regra de Hebb.} Dados $K$ padrões $\xi^1,\dots,\xi^K\in\{-1,+1\}^N$,
\begin{equation}\label{eq:hebb}
w_{ij}=\frac1N\sum_{\mu=1}^{K}\xi_i^{\mu}\xi_j^{\mu}\quad(i\neq j),\qquad w_{ii}=0,\qquad\theta=0 .
\end{equation}
Neurônios que têm o mesmo valor nos padrões ganham peso positivo, e os de valores opostos, peso
negativo~\cite{wikihopfield}. A regra é \emph{local} (o peso $w_{ij}$ depende apenas da atividade de $i$ e
$j$) e \emph{incremental} (um novo padrão é somado sem revisitar os anteriores). Hopfield e
Tank~\cite[p.~145]{hopfieldtank1985} chegam a $T_{ij}=\sum_sV_i^sV_j^s$, igual a~\eqref{eq:hebb} a menos
de um fator positivo e de uma constante aditiva na energia.\footnote{A Wikipédia~\cite{wikihopfield}
normaliza pelo número de padrões. Como o fator positivo não altera o sinal de $h_i$ quando $\theta=0$, a
escolha é irrelevante para a dinâmica; usamos $1/N$ para que a energia se escreva em função das
sobreposições (Apêndice~\ref{apx:deriv}).}

Aplicando a rede a um padrão armazenado $\xi^\nu$, o campo local se decompõe em sinal e ruído cruzado:
\[
h_i(\xi^\nu)=\underbrace{\tfrac{N-1}{N}\,\xi_i^\nu}_{\text{sinal}}+\underbrace{\tfrac1N\sum_{\mu\neq\nu}\sum_{j\neq i}\xi_i^\mu\xi_j^\mu\xi_j^\nu}_{\text{ruído cruzado}} .
\]
Para padrões independentes e não viesados, o ruído cruzado tem média nula e variância ${\approx}(K-1)/N$:
enquanto $K/N$ é pequeno, o sinal domina e o padrão é estável (argumento padrão; ver
Hertz et al.~\cite{hertz1991}). Daí a importância da razão $K/N$ na Seção~\ref{sec:capacidade}.

\paragraph{Regra de Storkey.} Introduzida em 1997~\cite{storkey1997}, também é local e incremental e, segundo o
autor, confere à rede capacidade maior que a da regra de Hebb~\cite{wikihopfield}. Com $W^{0}=0$ e $N$ neurônios,
\begin{equation}\label{eq:storkey}
w_{ij}^{\nu}=w_{ij}^{\nu-1}+\tfrac1N\xi_i^\nu\xi_j^\nu-\tfrac1N\xi_i^\nu h_{ji}^{\nu}-\tfrac1N\xi_j^\nu h_{ij}^{\nu},
\qquad h_{ij}^{\nu}=\sum_{k\neq i,j}w_{ik}^{\nu-1}\xi_k^\nu ,
\end{equation}
em que $h_{ij}^\nu$ é um campo local que exclui as contribuições de $i$ e $j$. \emph{Interpretação nossa:} os
termos de campo local subtraem do incremento a parte que os pesos já acumulados ``predizem'' para o novo padrão,
o que reduz o ruído cruzado. Em forma matricial, com $u=W^{\nu-1}\xi^\nu$ e a diagonal zerada após cada passo,
\[
W^{\nu}=W^{\nu-1}+\tfrac1N\bigl(\xi^\nu\xi^{\nu\top}-\xi^\nu u^{\top}-u\,\xi^{\nu\top}+2W^{\nu-1}\bigr),
\]
forma que verificamos numericamente contra a recursão~\eqref{eq:storkey} (Apêndice~\ref{apx:algoritmos}).

\subsection{Capacidade, estados espúrios e bacias de atração}\label{sec:capacidade}
\paragraph{Capacidade.} No espaço de dimensão $N$, o modelo padrão armazena apenas ${\sim}CN/\log N$ padrões
aleatórios sem erros, com $C<1/2$ para que um padrão fixo seja estável e $C<1/4$ para que todos o sejam
(McEliece et al.~\cite{mceliece1987}, conforme relatado em~\cite{ramsauer2021}). Aceitando pequenos erros de
recuperação, a capacidade é de cerca de $0{,}138N$~\cite{hertz1991,wikihopfield,ramsauer2021}. A Wikipédia
dá ainda $C\simeq N/(2\ln N)$ para a regra de Hebb e $C\simeq N/\sqrt{2\ln N}$ para a de
Storkey~\cite{wikihopfield}. A Tabela~\ref{tab:capacidade} avalia essas expressões para os tamanhos de rede
que surgem na Seção~\ref{sec:codificacao}.

\begin{table}[htbp]
\centering
\caption{Capacidades teóricas aproximadas (padrões aleatórios, independentes e não viesados). Valores
calculados pelas fórmulas citadas; são ordens de grandeza, não garantias, para $N$ desta faixa.}
\label{tab:capacidade}
\begin{tabular}{rccc}
\toprule
$N$ & $0{,}138\,N$ (Hebb, com pequenos erros) & $N/(2\ln N)$ (Hebb) & $N/\sqrt{2\ln N}$ (Storkey)\\
\midrule
67  & 9,2  & 8,0  & 23,1\\
122 & 16,8 & 12,7 & 39,4\\
224 & 30,9 & 20,7 & 68,1\\
412 & 56,9 & 34,2 & 118,7\\
\bottomrule
\end{tabular}
\end{table}

\paragraph{Estados espúrios.} Os padrões armazenados são atratores, mas a rede também pode convergir para
estados que não são memórias~\cite{wikihopfield}. Com $\theta=0$, $E(s)=E(-s)$, de modo que $-\xi^\mu$ é
sempre um estado espúrio com a mesma energia de $\xi^\mu$; combinações lineares de um número ímpar de
padrões também podem ser pontos fixos~\cite{wikihopfield}. O número de mínimos espúrios é exponencial na
dimensão~\cite{ramsauer2021,bruckroychowdhury1990}, mesmo para padrões ortogonais~\cite{wikihopfield}. Para
um detector de anomalias isso é um risco concreto: uma observação anômala atraída por um estado espúrio
termina com energia baixa e pode parecer familiar (Apêndice~\ref{apx:exemplo}).

\paragraph{Bacias de atração.} A bacia de um atrator é o conjunto de estados iniciais que convergem para
ele. Seu tamanho determina a tolerância a ruído da recuperação; no nosso caso, define o quanto uma
observação normal pode variar sem ser tratada como anomalia.

\subsection{Rede contínua e a formulação de Hopfield e Tank}\label{sec:ht}
Hopfield e Tank~\cite{hopfieldtank1985} usam neurônios de resposta graduada, $V_i=g(u_i)$ com $g$ sigmoide, e
dinâmica
\begin{equation}\label{eq:ht-dyn}
\frac{du_i}{dt}=-\frac{u_i}{\tau}+\sum_jT_{ij}V_j+I_i\qquad\text{(Eq.~3 do artigo)},
\end{equation}
em que $T_{ij}$ são as conexões e $I_i$ correntes de polarização. Com $T$ simétrica o sistema converge a estados estáveis~\cite{hopfield1984,hopfieldtank1985}, que no limite de alto ganho são os mínimos locais de
\begin{equation}\label{eq:ht-energia}
E=-\tfrac12\sum_i\sum_jT_{ij}V_iV_j-\sum_iV_iI_i\qquad\text{(Eq.~4)},
\end{equation}
situados nos vértices do hipercubo~\cite{hopfieldtank1985}.

\paragraph{A CAM como problema de otimização.} Na CAM (pp.~144--145) as saídas valem $\pm1$ e a memória é uma
palavra binária; a recuperação coloca as saídas no padrão da chave e deixa o sistema convergir. Para que os
mínimos de $E$ sejam as memórias $V^s$, os autores propõem $E=-\tfrac12\sum_s(V^s\cdot V)^2$ (Eq.~5), que
tem mínimos de profundidade ${\approx}-\tfrac12N^2$ em cada memória, e mostram que isso equivale
a~\eqref{eq:ht-energia} com $I_i=0$ e $T_{ij}=\sum_sV_i^sV_j^s$. Na p.~145 observam também que, numa CAM
prática (por exemplo, para guardar números de telefone), é preciso uma transformação que codifique a
informação do mundo real na representação binária da memória. Esse é o papel da codificação em termômetro
da Seção~\ref{sec:codificacao}.

\paragraph{O caixeiro-viajante.} Os autores mapeiam o problema em $n^2$ neurônios (a posição de cada cidade no
percurso), com termos de energia para as restrições e para a distância. Com 10 cidades, de 20 estados
iniciais aleatórios 16 convergiram para percursos válidos, e cerca de 50\,\% das tentativas produziram um dos
2 menores percursos; com 30 cidades (900 neurônios) a escolha de parâmetros ``parece ser uma questão mais
delicada'' do que com 100 neurônios~\cite[p.~148]{hopfieldtank1985}. Para este trabalho, a lição é dupla:
mapear o problema na rede é parte central do modelo, e a qualidade do resultado depende de sintonia de
parâmetros que se agrava com o tamanho. Não usamos a dinâmica analógica nem otimização combinatória; usamos a
rede apenas como memória associativa.

\subsection{Redes de Hopfield modernas}\label{sec:moderna}
Krotov e Hopfield~\cite{krotov2016} trocaram o termo quadrático da energia por uma função de crescimento mais
rápido ($F(x)=x^n$), obtendo capacidade proporcional a $N^{n-1}$; Demircigil et al.~\cite{demircigil2017}
usaram energia exponencial, com capacidade exponencial em $N$~\cite{wikihopfield}. Ramsauer
et al.~\cite{ramsauer2021} generalizam para estados contínuos $\xi\in\R^N$ e padrões nas colunas de
$X=(x_1,\dots,x_K)\in\R^{N\times K}$, com energia e regra de atualização
\begin{align}
E(\xi)&=-\lse(\beta,X^{\top}\xi)+\tfrac12\xi^{\top}\xi+\beta^{-1}\ln K+\tfrac12M^2,\qquad M=\max_i\|x_i\|,\label{eq:energia-moderna}\\
\xi^{\text{novo}}&=X\,\softmax\bigl(\beta X^{\top}\xi\bigr),\label{eq:update-moderna}
\end{align}
em que $\lse(\beta,z)=\beta^{-1}\ln\sum_i\exp(\beta z_i)$. Os resultados principais do artigo são: a iteração
converge globalmente (Teoremas 1 e 2); para padrões aleatórios na esfera, com probabilidade $1-p$ é possível
armazenar ao menos $\sqrt p\,c^{(N-1)/4}$ padrões (Teorema~3), em que $c$ depende de $\beta$, do raio da esfera,
da dimensão e de $p$ --- o artigo demonstra $c\ge3{,}1546$ e $c\ge1{,}3718$ para dois conjuntos de parâmetros,
de modo que o limite cresce exponencialmente com $N$; e a recuperação ocorre tipicamente em uma atualização, com
erro exponencialmente pequeno na separação $\Delta_i=x_i^{\top}x_i-\max_{j\neq i}x_i^{\top}x_j$
(Teoremas 4 e 5). O parâmetro $\beta$ controla o tipo de
ponto fixo: média global de todos os padrões, estado metaestável (média de um subconjunto de padrões
semelhantes) ou recuperação de um único padrão. A regra~\eqref{eq:update-moderna} coincide com a atenção do
\emph{Transformer}, o que levou os autores a propor camadas Hopfield (\texttt{Hopfield},
\texttt{HopfieldPooling} e \texttt{HopfieldLayer})~\cite{ramsauer2021,hopfieldlayers}.

\paragraph{Interpretação (derivação própria).} Substituindo $x_i^{\top}\xi=\tfrac12\|x_i\|^2+\tfrac12\|\xi\|^2-\tfrac12\|\xi-x_i\|^2$
em~\eqref{eq:energia-moderna}, o termo $\tfrac12\xi^\top\xi$ se cancela e resulta (Apêndice~\ref{apx:deriv})
\[
E(\xi)=-\beta^{-1}\ln\sum_{i=1}^{K}e^{\frac12\beta\|x_i\|^2}e^{-\frac12\beta\|\xi-x_i\|^2}+\beta^{-1}\ln K+\tfrac12M^2 .
\]
Se todos os padrões têm a mesma norma, $E(\xi)=-\beta^{-1}\ln\bigl[\tfrac1K\sum_ie^{-\beta\cdot\frac12\|\xi-x_i\|^2}\bigr]$,
um \emph{mínimo suave} (log-média-exp) das meias distâncias quadráticas aos padrões. Ele tende a
$\min_i\tfrac12\|\xi-x_i\|^2$ (vizinho mais próximo) quando $\beta\to\infty$ e à média dessas distâncias
(cujo mínimo é o centroide) quando $\beta\to0$. Além disso, $e^{-\beta E(\xi)}$ é, a menos de constante,
uma estimativa de densidade por kernel gaussiano com variância $\sigma^2=1/\beta$.
\textbf{Consequência para este trabalho:} um escore de anomalia baseado em $E(\xi)$ é essencialmente um escore
de novidade do tipo KDE/\emph{k}-NN. Não esperamos grande vantagem sobre essas linhas de base; o valor da rede
moderna está na recuperação associativa iterativa e na ligação com a atenção.

\begin{table}[htbp]
\centering
\caption{Rede de Hopfield clássica e moderna.}
\label{tab:comparacao}
\small
\begin{tabularx}{\textwidth}{@{}l L L@{}}
\toprule
Aspecto & Clássica (Hebb/Storkey) & Moderna (Ramsauer et al.)\\
\midrule
Estado & $s\in\{-1,+1\}^N$ & $\xi\in\R^N$ (contínuo)\\
Memória & $W\in\R^{N\times N}$: $O(N^2)$ valores, independente de $K$ & $X\in\R^{N\times K}$: $O(NK)$ valores (guarda os padrões)\\
Energia & $-\tfrac12s^\top Ws$ & $-\lse(\beta,X^\top\xi)+\tfrac12\xi^\top\xi+\cdots$\\
Atualização & $s_i\leftarrow\sgn(\sum_jw_{ij}s_j)$, assíncrona & $\xi\leftarrow X\softmax(\beta X^\top\xi)$\\
Capacidade (padrões aleatórios) & ${\approx}0{,}138N$ (Hebb); ${\approx}N/\sqrt{2\ln N}$ (Storkey) & exponencial em $N$ (Teorema~3)\\
Recuperação & iterativa (varreduras) & tipicamente uma atualização (Teorema~4)\\
Pontos fixos & padrões, $-$padrões e estados espúrios & um padrão, estados metaestáveis ou média global (via $\beta$)\\
Hiperparâmetros & regra de armazenamento, $\theta$ & $\beta$\\
Fontes & \cite{hopfield1982,storkey1997,wikihopfield} & \cite{ramsauer2021,krotov2016,demircigil2017}\\
\bottomrule
\end{tabularx}
\end{table}

\subsection{A rede de Hopfield como detector de anomalias}\label{sec:conceito}
Seja $\mathcal M=\{\xi^1,\dots,\xi^K\}$ o conjunto de padrões normais gravados e $q$ uma observação,
codificada como estado inicial $s^0$. A rede atua como \emph{filtro de recuperação}, $s^0\mapsto\sstar$.
Esperamos três regimes (Figura~\ref{fig:energia}):
\begin{enumerate}
  \item \textbf{$q$ normal:} $s^0$ cai na bacia de um protótipo; $\sstar\approx\xi^\mu$, $E(s^0)$ é baixa e
        poucos bits mudam.
  \item \textbf{$q$ anômalo fora das bacias:} $E(s^0)$ é alta; a dinâmica altera muitos bits ou termina longe
        de todos os protótipos.
  \item \textbf{$q$ anômalo capturado por atrator espúrio} (mistura de padrões ou padrão invertido
        $-\xi^\mu$): $\sstar$ tem energia baixa mas não está próximo de nenhum protótipo; só escores que medem
        a distância a $\mathcal M$ o denunciam (Seção~\ref{sec:escores}).
\end{enumerate}
O treinamento é \emph{one-class}: usa somente observações predominantemente normais e nenhum rótulo. Os
rótulos entram apenas para escolher o limiar de decisão e para avaliar o detector.

\begin{figure}[htbp]
\centering
\newcommand{\Efun}[1]{-2*exp(-((#1-1.8)^2)/0.5)-2.6*exp(-((#1-5)^2)/0.5)-2*exp(-((#1-8)^2)/0.5)-0.7*exp(-((#1-6.6)^2)/0.25)}
\begin{tikzpicture}[x=1.35cm,y=1.1cm,font=\small]
  \draw[->] (0,-3.2) -- (10.4,-3.2) node[right] {estados};
  \draw[->] (0,-3.2) -- (0,0.7) node[above] {$E$};
  \draw[very thick,smooth,domain=0.2:9.8,samples=240,variable=\x]
    plot ({\x},{\Efun{\x}});
  % protótipos
  \foreach \c/\n in {1.8/1,5/2,8/3}{
    \pgfmathsetmacro{\yc}{\Efun{\c}}
    \fill (\c,\yc) circle (2pt);
    \node[below=2pt] at (\c,\yc) {$\xi^{\n}$};
  }
  % atrator espúrio
  \pgfmathsetmacro{\ys}{\Efun{6.6}}
  \draw[red!70!black,thick] (6.6,\ys) circle (2.5pt);
  \node[red!70!black,below=3pt,align=center,font=\scriptsize] at (6.6,\ys) {atrator\\espúrio};
  % observação normal desce para o protótipo
  \pgfmathsetmacro{\yn}{\Efun{2.7}}
  \fill[blue!70!black] (2.7,\yn) circle (2.2pt);
  \draw[-{Stealth[length=2mm]},blue!70!black,thick,smooth,domain=2.6:2.05,samples=20,variable=\x]
    plot ({\x},{\Efun{\x}});
  \node[blue!70!black,above left=1pt and 2pt,align=right,font=\scriptsize] at (2.7,\yn)
    {observação\\normal};
  % observação anômala no platô de energia alta
  \pgfmathsetmacro{\ya}{\Efun{3.6}}
  \fill[red!70!black] (3.6,\ya) circle (2.2pt);
  \node[red!70!black,above=3pt,align=center,font=\scriptsize] at (3.95,\ya)
    {observação anômala\\(energia alta)};
  % observação anômala capturada por estado espúrio
  \pgfmathsetmacro{\yb}{\Efun{6.3}}
  \fill[red!70!black] (6.3,\yb) circle (2.2pt);
  \draw[-{Stealth[length=2mm]},red!70!black,thick,dashed,smooth,domain=6.3:6.55,samples=10,variable=\x]
    plot ({\x},{\Efun{\x}});
\end{tikzpicture}
\caption{Paisagem de energia esquemática, unidimensional: três protótipos normais, um atrator espúrio e três
observações. Ilustração conceitual; os valores são arbitrários e não representam resultados.}
\label{fig:energia}
\end{figure}

\section{Descrição do problema e dos dados}\label{sec:dados}

\subsection{Definição do problema}\label{sec:problema}
Considere uma máquina monitorada continuamente. A cada minuto $t$ são coletadas $D=38$ métricas de utilização
e de atividade, que formam o vetor $x_t\in[0,1]^{D}$. Especialistas de domínio marcaram os intervalos em que a
máquina esteve em incidente, o que fornece o rótulo $y_t\in\{0,1\}$ ($0$ = normal, $1$ = anômalo). A tarefa é
uma \emph{classificação binária por ponto}: dado $x_t$ (opcionalmente com uma janela de observações anteriores),
prever $y_t$ por
\begin{equation}\label{eq:decisao}
\hat y_t=\mathbf 1\bigl[\,a(x_t)>\tau\,\bigr],
\end{equation}
em que $a(\cdot)$ é um \emph{escore de anomalia} (quanto maior, mais anômalo) e $\tau$ é um limiar escolhido em
dados de validação. A rede de Hopfield produz o escore $a$ (Seção~\ref{sec:escores}); as linhas de base
produzem $\hat y_t$ diretamente ou por um escore equivalente.

\begin{itemize}
  \item \textbf{Entrada (atributos):} $x_t\in[0,1]^{38}$, com 38 métricas numéricas anônimas (``carga de CPU,
        uso de rede, uso de memória etc.''~\cite{su2019}), sem nomes nem unidades, amostradas a cada minuto.
  \item \textbf{Saída desejada:} o rótulo $y_t\in\{0,1\}$ e, no caso do detector, o escore $a(x_t)$.
  \item \textbf{Unidade de modelagem:} cada máquina é tratada separadamente, com seus próprios dados de
        treino e de teste~\cite{su2019}.
\end{itemize}

\subsection{O conjunto de dados SMD}\label{sec:smd}
Adotamos o \emph{Server Machine Dataset} (SMD), publicado por Su et al.~\cite{su2019} e disponibilizado no
repositório do OmniAnomaly~\cite{omnianomaly}. Ele reúne cerca de cinco semanas de telemetria de 28 máquinas
(``entidades'') de uma grande empresa de Internet, divididas em três grupos de 8, 9 e 11 máquinas, com 38
métricas por máquina e amostragem de 1 minuto. Para cada máquina, a primeira metade da série é o conjunto de
treino (\emph{sem rótulos}) e a segunda é o conjunto de teste, rotulado por especialistas de domínio a partir de
relatórios de incidentes~\cite{su2019}. Os dados \emph{não} são redistribuídos neste projeto
(Seção~\ref{sec:repro}).

Calculamos as estatísticas a seguir diretamente dos arquivos do repositório. Há 708.405 linhas de treino e
708.420 de teste (as duas metades de uma máquina diferem em no máximo uma linha). Dos pontos de teste,
29.444 (4,16\,\%) são anômalos e se agrupam em 327 segmentos contíguos, com duração média de 90,0 min e
extremos de 2 e 3.161 min. Um classificador que responde sempre ``normal'' acerta 95,84\,\% dos pontos de
teste; portanto a acurácia, isoladamente, é uma métrica inadequada (Seção~\ref{sec:metricas}). A
Tabela~\ref{tab:grupos} resume os três grupos; o Apêndice~\ref{apx:maquinas} traz os números por máquina.

\begin{table}[htbp]
\centering
\caption{SMD por grupo. Estatísticas calculadas pelos autores sobre os arquivos do repositório. Anôm. =
pontos anômalos do teste; Seg. = segmentos anômalos; Prev. = prevalência de anomalias por máquina (mínimo e
máximo do grupo).}
\label{tab:grupos}
\small
\setlength{\tabcolsep}{5pt}
\begin{tabular}{crrrrcr}
\toprule
Grupo & Máquinas & Pontos de teste & Anôm. (\%) & Seg. & Prev. por máquina & Linhas de treino\\
\midrule
1 & 8  & 194.374 & 11.742 (6,04) & 112 & 0,42--15,65\,\% & 194.369\\
2 & 9  & 223.325 &  9.703 (4,34) & 114 & 0,68--11,95\,\% & 223.320\\
3 & 11 & 290.721 &  7.999 (2,75) & 101 & 0,69--4,78\,\%  & 290.716\\
\midrule
Total & 28 & 708.420 & 29.444 (4,16) & 327 & 0,42--15,65\,\% & 708.405\\
\bottomrule
\end{tabular}
\end{table}

Na máquina 1-1, que usamos como exemplo ao longo do texto, cada metade tem $28.479\times38$ valores, todos em
$[0,1]$, sem valores ausentes e sem linhas de treino duplicadas. A verificação das demais máquinas faz parte
da preparação (Seção~\ref{sec:prep}).

\subsection{Três exemplos reais}\label{sec:exemplos}
A Tabela~\ref{tab:exemplos} mostra três observações reais da máquina \texttt{machine-1-1}: uma normal
(linha 15.000 do arquivo de teste) e duas anômalas (linhas 16.100 e 17.200, dentro dos segmentos 15.850--16.395
e 16.964--17.517). Por brevidade, mostramos 10 dos 38 atributos --- os que o arquivo
\texttt{interpretation\_label} associa a esses segmentos --- junto com a mediana, o percentil 95 e o máximo
do \emph{treino}. O arquivo \texttt{interpretation\_label} indica, para cada segmento anômalo, os atributos
que o explicam (numerados a partir de 1); seus intervalos diferem ligeiramente dos do arquivo de rótulos
(por exemplo, 15.849--16.368 contra 15.850--16.395).

\begin{table}[htbp]
\centering
\caption{Três exemplos reais da máquina \texttt{machine-1-1} (linhas do arquivo de teste, contadas a partir de
1) e estatísticas dos mesmos atributos no treino. Atributos numerados a partir de 1. Em negrito, valor acima do
máximo do treino. O atributo 29 é constante (igual a zero) em todo o treino.}
\label{tab:exemplos}
\small
\setlength{\tabcolsep}{5pt}
\begin{tabular}{rcccccc}
\toprule
 & \multicolumn{3}{c}{Teste} & \multicolumn{3}{c}{Treino}\\
\cmidrule(lr){2-4}\cmidrule(lr){5-7}
Atributo & linha 15.000 & linha 16.100 & linha 17.200 & mediana & p95 & máximo\\
 & ($y=0$) & ($y=1$) & ($y=1$) & & &\\
\midrule
1  & 0,0968 & \textbf{0,5161} & \textbf{0,5161} & 0,0430 & 0,1398 & 0,4946\\
2  & 0,0794 & 0,0816 & 0,1250 & 0,0456 & 0,1409 & 0,4216\\
9  & 0,0288 & 0,0129 & 0,0254 & 0,0210 & 0,0355 & 0,1181\\
10 & 0,0000 & 0,0035 & 0,0182 & 0,0000 & 0,0004 & 0,1404\\
12 & 0,1081 & 0,2297 & 0,1351 & 0,0946 & 0,1622 & 0,7432\\
13 & 0,0274 & 0,0548 & 0,0548 & 0,0274 & 0,0548 & 0,3151\\
14 & 0,0575 & 0,0314 & 0,0700 & 0,0500 & 0,0781 & 0,1231\\
15 & 0,1250 & 0,0684 & 0,0965 & 0,0969 & 0,1395 & 0,5650\\
25 & 0,0493 & 0,0313 & 0,0364 & 0,0145 & 0,0561 & 0,1796\\
29 & 0,0000 & 0,0000 & 0,0000 & 0,0000 & 0,0000 & 0,0000\\
\bottomrule
\end{tabular}
\end{table}

No exemplo da linha 16.100, o atributo 1 sobe a 0,5161, acima de todo o treino (máximo 0,4946), e o atributo 12
chega a 0,2297, acima do percentil 95 do treino (0,1622). No exemplo da linha 17.200, o atributo 1 repete o
valor 0,5161, mas os demais desvios são pequenos em valor absoluto: o atributo 10 vale 0,0182, muito acima do
percentil 95 do treino (0,0004) e ainda assim dentro do intervalo observado (máximo 0,1404). A tarefa já
mostra, portanto, uma dificuldade típica: ``anômalo'' nem sempre significa ``fora da faixa''. Em toda a
máquina 1-1, apenas 18,19\,\% dos pontos anômalos têm o atributo 1 acima do máximo do treino (contra 0,19\,\%
dos normais) e apenas 1,00\,\% têm o atributo 10 acima do seu máximo (contra 0,00\,\% dos normais).

\subsection{Exploração dos dados: atributos constantes e deriva}\label{sec:exploracao}
\paragraph{Atributos constantes.} Na máquina 1-1, 8 dos 38 atributos são constantes no treino (atributos 5, 8,
17, 18, 27, 29, 37 e 38), restando 30 atributos variáveis. No teste, o atributo 27 deixa de ser constante em
22 linhas (todas anômalas) e o atributo 29 em 225 linhas (211 anômalas e 14 normais). Um atributo que ``nunca
se mexia'' e passa a se mexer é um sinal de novidade que a codificação precisa preservar
(Seção~\ref{sec:codificacao}).

\paragraph{Deriva entre treino e teste.} Como o teste vem depois do treino, a distribuição dos atributos pode
mudar mesmo sem anomalias (\emph{drift}). Para quatro máquinas (1-1, 1-3, 2-5 e 3-6) comparamos, atributo a
atributo, a distribuição do treino com a das linhas \emph{normais} do teste pela estatística de
Kolmogorov--Smirnov (KS), e contamos as linhas com ao menos um atributo acima do máximo do treino
(Tabela~\ref{tab:deriva}).

\begin{table}[htbp]
\centering
\caption{Deriva entre treino e teste em quatro máquinas (cálculo próprio). KS: estatística de Kolmogorov--Smirnov
entre o treino e as linhas normais do teste, sobre os atributos variáveis. ``Acima do máx.'': percentual de
linhas do teste com ao menos um atributo acima do máximo do treino.}
\label{tab:deriva}
\small
\setlength{\tabcolsep}{4pt}
\begin{tabularx}{\textwidth}{@{}l>{\hsize=.75\hsize}Y>{\hsize=.85\hsize}Y>{\hsize=.7\hsize}Y>{\hsize=1.2\hsize}Y>{\hsize=1.5\hsize}L@{}}
\toprule
Máquina & Atributos variáveis & Atributos com KS\,$>0{,}2$ & KS mediana & Acima do máx.: normais / anômalas & Atributos constantes no treino\\
\midrule
1-1 & 30 & 22 & 0,37 & 30,1\,\% / 71,6\,\% & 8 (5, 8, 17, 18, 27, 29, 37, 38)\\
1-3 & 34 & 7  & 0,08 & 0,9\,\% / 18,0\,\%  & 4 (8, 27, 37, 38)\\
2-5 & 31 & 2  & 0,06 & 1,1\,\% / 17,2\,\%  & 7 (5, 8, 17, 18, 29, 37, 38)\\
3-6 & 30 & 7  & 0,13 & 6,5\,\% / 35,8\,\%  & 8 (5, 8, 17, 18, 27, 29, 37, 38)\\
\bottomrule
\end{tabularx}
\end{table}

A deriva é \emph{heterogênea}: forte na máquina 1-1 (22 dos 30 atributos com KS\,$>0{,}2$; KS máximo de 0,70 no
atributo 30) e leve em 1-3 e 2-5. Na máquina 1-1, o atributo 30 sozinho responde por 26,79\,\% das linhas normais
acima do máximo do treino (atributos 4 e 3: 5,41\,\% e 3,91\,\%), e apenas 0,02\,\% das linhas têm algum atributo
abaixo do mínimo do treino. Excluindo esses três atributos, o percentual de linhas acima do máximo cai para
5,4\,\% nas normais contra 43,5\,\% nas anômalas. Duas conclusões orientam a metodologia: (i) ``acima do máximo
do treino'' \emph{sozinho} é um detector pobre, pois dispara em muitas linhas normais quando há deriva e deixa
passar muitas anomalias quando não há; (ii) a calibração do limiar e, em implantações reais, a atualização
periódica do modelo são essenciais.

\subsection{Adequação ao tema e limitações do conjunto}\label{sec:adequacao}
\paragraph{Pontos favoráveis.} Os dados são reais, de produção, com métricas de utilização e de atividade
por máquina; os rótulos vêm de incidentes registrados; e o treino é \emph{sem rótulos}, o que combina com o
treinamento \emph{one-class} da rede de Hopfield (Seção~\ref{sec:conceito}).

\paragraph{Limitações.}
\begin{itemize}
  \item \textbf{``Baixo nível'' é aproximado.} São 38 contadores anonimizados de utilização de recursos, sem
        chamadas de sistema nem rastros de \emph{kernel}. O SMD é um substituto para estudar a abordagem; a
        coleta própria de telemetria (\texttt{psutil}, eBPF) com injeção de falhas é trabalho
        futuro (Seção~\ref{sec:conclusoes}).
  \item \textbf{Atributos anônimos.} Sem nomes ou unidades, não há semântica a explorar na interpretação dos
        resultados nem como conhecimento prévio de um LLM (Etapa~II).
  \item \textbf{Rótulos de incidentes.} As fronteiras dos segmentos são imprecisas (basta ver a diferença entre
        os arquivos de rótulos e de interpretação) e o rótulo vale para o intervalo inteiro, mesmo que o
        efeito nas métricas varie dentro dele.
  \item \textbf{Treino sem rótulos.} A metade de treino pode conter anomalias não marcadas; assumimos que é
        majoritariamente normal.
  \item \textbf{Heterogeneidade.} A prevalência varia de 0,42\,\% a 15,65\,\% por máquina, e a deriva também
        (Tabela~\ref{tab:deriva}). Várias máquinas não têm anomalias em algum dos blocos cronológicos usados
        para treinar, validar e testar (Seção~\ref{sec:divisao}).
  \item \textbf{Licença.} O repositório do OmniAnomaly declara licença MIT; confirmar os termos de uso dos
        dados antes de qualquer redistribuição.
\end{itemize}

\section{Metodologia}\label{sec:metodo}

\subsection{Visão geral}\label{sec:visao}
A Figura~\ref{fig:pipeline} resume a abordagem. Na fase \emph{offline}, a metade de treino (sem rótulos) é
limpa, binarizada por uma codificação em termômetro, resumida em $K$ protótipos, e esses protótipos são gravados
na matriz de pesos $W$. Na fase \emph{online}, cada observação é codificada, a rede a leva a um estado estável e
a diferença entre o estado inicial e o recuperado (ou a energia) vira um escore de anomalia, que um limiar
converte em $\hat y_t$. Rótulos só são usados para escolher o limiar e para avaliar (Seção~\ref{sec:divisao}).

\begin{figure}[htbp]
\centering
\begin{tikzpicture}[node distance=6mm,
  box/.style={draw,rounded corners=2pt,thick,minimum width=25mm,minimum height=9mm,align=center,
              font=\footnotesize,fill=blue!6},
  mem/.style={box,fill=orange!18},
  arr/.style={-{Stealth[length=2mm]},thick},
  dotarr/.style={-{Stealth[length=2mm]},thick,dashed,gray!70!black}]
  \node[box] (pre) {limpeza e\\preparação};
  \node[box,right=of pre] (cod) {codificação\\em termômetro};
  \node[box,right=of cod] (pro) {protótipos\\($K$ padrões)};
  \node[mem,right=of pro] (mem) {memória $W$\\(Hebb/Storkey)};
  \draw[arr] (pre)--(cod);
  \draw[arr] (cod)--(pro);
  \draw[arr] (pro)--(mem);
  \node[anchor=south west,font=\footnotesize\bfseries,text=blue!50!black]
    at ($(pre.north west)+(0,2mm)$) {Fase offline (treino, sem rótulos)};

  \node[box,below=14mm of pre] (obs) {observação\\$x_t$};
  \node[box,right=of obs] (cod2) {codificação\\$s^{0}$};
  \node[box,right=of cod2] (rec) {recuperação\\assíncrona $\sstar$};
  \node[box,right=of rec] (esc) {escore\\$a(x_t)$};
  \node[box,right=of esc,minimum width=17mm] (dec) {$a>\tau$?\\$\hat y_t$};
  \draw[arr] (obs)--(cod2);
  \draw[arr] (cod2)--(rec);
  \draw[arr] (rec)--(esc);
  \draw[arr] (esc)--(dec);
  \node[anchor=north west,font=\footnotesize\bfseries,text=blue!50!black]
    at ($(obs.south west)+(0,-2mm)$) {Fase online (inferência)};

  \draw[dotarr] (cod.south)--(cod2.north);
  \draw[dotarr] (mem.south)--++(0,-5mm)-|(rec.north);
\end{tikzpicture}
\caption{Pipeline do detector. Setas tracejadas: parâmetros aprendidos \emph{offline} (limiares da codificação
e matriz de pesos) e reutilizados \emph{online}.}
\label{fig:pipeline}
\end{figure}

\subsection{Preparação e limpeza}\label{sec:prep}
\begin{enumerate}
  \item \textbf{Verificação por máquina:} dimensões ($n\times38$), valores ausentes, valores fora de $[0,1]$,
        linhas duplicadas e coerência entre o número de linhas do teste e o do arquivo de rótulos. Na máquina
        1-1 não há problemas (Seção~\ref{sec:smd}). Se houver valores ausentes em outras máquinas, serão
        preenchidos com o último valor válido, nunca com estatísticas do teste.
  \item \textbf{Ordem temporal preservada:} nada é embaralhado.
  \item \textbf{Sem normalização adicional:} os valores já estão em $[0,1]$. Os modelos que exigem
        padronização (SVM, MLP, \emph{k}-NN, OC-SVM) usam média e desvio-padrão calculados \emph{apenas} no
        treino; atributos constantes recebem escala~1, de modo que um desvio no teste continue visível.
  \item \textbf{Atributos constantes} no treino são mantidos e tratados por um bit de novidade
        (Seção~\ref{sec:codificacao}).
  \item \textbf{Treino sem remoção de \emph{outliers}:} sem rótulos, não há como identificar com segurança as
        anomalias eventualmente presentes no treino; a mitigação é a seleção de protótipos com descarte de
        grupos muito pequenos (Seção~\ref{sec:prototipos}).
  \item \textbf{Janelas (ablação):} opcionalmente, $x_t$ é ampliado com resumos (média e máximo por atributo) dos
        últimos $w$ minutos, para capturar a dinâmica temporal. O modelo-base usa só a observação instantânea.
\end{enumerate}

\subsection{Codificação binária em termômetro}\label{sec:codificacao}
A rede de Hopfield discreta trabalha com estados em $\{-1,+1\}^N$, enquanto $x_t\in[0,1]^{38}$. Hopfield e
Tank~\cite[p.~145]{hopfieldtank1985} observam que, numa memória prática, é preciso uma transformação que
codifique a informação do mundo real na representação binária da memória. Usamos a codificação em
\emph{termômetro}.

Para cada atributo variável $j$ e um número de níveis $b$, os limiares $\tau_{j,l}$ são os quantis empíricos do
treino de ordem $l/(b+1)$, $l=1,\dots,b$, descartando repetidos e os que não satisfazem
$\min_t x_{t,j}\le\tau<\max_t x_{t,j}$ (assim cada bit varia no treino). O código do atributo é
$s_{j,l}=+1$ se $x_j>\tau_{j,l}$ e $s_{j,l}=-1$ caso contrário. Cada atributo \emph{constante} no treino, com
valor $c_j$, recebe um único \emph{bit de novidade}: $+1$ se $|x_j-c_j|>\varepsilon$ ($\varepsilon=10^{-6}$) e
$-1$ caso contrário. O estado total tem $N=N_{\text{term}}+n_{\text{const}}$ bits.

Uma propriedade útil (derivada no Apêndice~\ref{apx:deriv}): para um mesmo atributo, a distância de Hamming
entre os códigos de dois valores é o número de limiares entre eles, isto é, a diferença de postos
(quantis). A distância de Hamming entre estados completos é, portanto, uma distância \emph{ordinal}, somada
entre atributos, o que a torna insensível à escala de cada atributo.

\begin{table}[htbp]
\centering
\caption{Exploração da codificação em termômetro na máquina 1-1 (metade de treino, 28.479 linhas; cálculo
próprio). ``Padrões distintos'': linhas de treino com códigos diferentes. $\bar s_k$: média do bit $k$ no
treino. A última coluna é a razão entre o número de linhas de treino e a capacidade clássica $0{,}138N$.}
\label{tab:codificacao}
\small
\setlength{\tabcolsep}{4pt}
\begin{tabularx}{\textwidth}{@{}rRRRRRR@{}}
\toprule
$b$ & Bits de termômetro & $N$ (com 8 de novidade) & Padrões distintos & Média dos bits & Média de $|\bar s_k|$ & Treino $/\,(0{,}138N)$\\
\midrule
2  & 59  & 67  & 22.855 & $-0{,}065$ & 0,333 & 3.080\\
4  & 114 & 122 & 27.670 & $-0{,}064$ & 0,405 & 1.692\\
8  & 216 & 224 & 28.459 & $-0{,}040$ & 0,451 & 921\\
16 & 404 & 412 & 28.479 & $-0{,}032$ & 0,480 & 501\\
\bottomrule
\end{tabularx}
\end{table}

A Tabela~\ref{tab:codificacao} traz dois alertas. Primeiro, o número de linhas de treino é de 501 a 3.080
vezes a capacidade clássica: \emph{armazenar todos os pontos é impossível}, e precisamos de $K\ll0{,}138N$
protótipos (ou da rede moderna, com memória grande). Segundo, os bits do termômetro são \emph{viesados por
construção}: o bit do quantil $l/(b+1)$ tem média $1-2l/(b+1)$ no treino, e a média de $|\bar s_k|$ fica entre
0,33 e 0,48. As estimativas de capacidade da Seção~\ref{sec:capacidade} supõem padrões não viesados, logo a
capacidade efetiva deve ser menor que a da Tabela~\ref{tab:capacidade}. Por fim, cada bit de novidade pesa
apenas $1/N$ em $S_1$ (Seção~\ref{sec:escores}), um sinal fraco; combinar o detector com uma regra explícita
é uma extensão possível (Seção~\ref{sec:conclusoes}).

\subsection{Memória de protótipos}\label{sec:prototipos}
Os $K$ padrões armazenados são escolhidos entre os códigos da metade de treino por \emph{k-medoids} no espaço
de Hamming (ou \emph{k-modes}), ajustado sobre uma subamostra de até 5.000 pontos para limitar o custo. Usar
medoides garante que cada protótipo seja um padrão realmente observado. Grupos com menos de uma pequena fração
dos pontos (por exemplo, 1\,\%) podem ser descartados, para diminuir o risco de gravar anomalias que estejam
no treino. Variamos $K\in\{2,5,10,20,30\}$; esse intervalo inclui valores acima da capacidade teórica (por
exemplo, $K=30$ com $N=67$, $K/N\approx0{,}45$) de propósito, para observar a degradação prevista (H1,
Seção~\ref{sec:experimentos}). Os protótipos são gravados com a regra de Hebb~\eqref{eq:hebb} ou de
Storkey~\eqref{eq:storkey}, com $\theta=0$ (Algoritmo~\ref{alg:armazenar}).

\subsection{Recuperação e escores de anomalia}\label{sec:escores}
Dada uma observação $x_t$, o detector (i) calcula $s^{0}=\mathrm{cod}(x_t)$; (ii) executa a recuperação
assíncrona (Algoritmo~\ref{alg:recuperar}): varreduras em ordem aleatória até que uma varredura completa não
altere nenhum bit, com no máximo $T_{\max}=100$ varreduras --- a Proposição~\ref{prop:conv} garante a parada ---
e obtém $\sstar$; (iii) calcula os escores da Tabela~\ref{tab:escores}.

\begin{table}[htbp]
\centering
\caption{Escores de anomalia da rede clássica (quanto maior, mais anômalo). $\xi^\mu$: protótipos
armazenados; $E$: energia~\eqref{eq:energia}. Propostas de projeto.}
\label{tab:escores}
\small
\begin{tabularx}{\textwidth}{@{}lcL@{}}
\toprule
Escore & Definição & Leitura\\
\midrule
$S_1$ & $\hamm(s^0,\sstar)/N$ & fração dos bits alterada pela dinâmica\\
$S_2$ & $\min_\mu\hamm(s^0,\xi^\mu)/N$ & distância ao protótipo mais próximo \emph{sem} dinâmica (referência do vizinho mais próximo no espaço do código)\\
$S_3$ & $E(s^0)$ & energia do estado inicial\\
$S_4$ & $E(s^0)-E(\sstar)$ & energia liberada pela recuperação\\
$S_5$ & $\min_\mu\hamm(\sstar,\xi^\mu)/N$ & distância do estado final ao protótipo mais próximo (alta $\Rightarrow$ estado espúrio ou invertido)\\
$S_6$ & nº de varreduras com ao menos uma troca & tempo de convergência\\
\bottomrule
\end{tabularx}
\end{table}

A comparação entre $S_1$ e $S_2$ mede se a dinâmica recorrente acrescenta algo à simples busca do protótipo
mais próximo (ablação). Usamos a distância \emph{com sinal}, e não o módulo da sobreposição, de modo que o padrão
invertido $-\xi^\mu$ --- que é ponto fixo com a mesma energia de $\xi^\mu$ quando $\theta=0$ (Seção~\ref{sec:capacidade})
--- conte como distante: para uma consulta igual a $-\xi^\mu$, $S_1=0$ e $S_3$ é baixo, mas $S_2=S_5=1$
(Apêndice~\ref{apx:exemplo}).

\paragraph{Variante moderna.} A memória é a matriz $X$ formada pelas linhas de treino (ou por uma subamostra
ou centroides de $k$-means), no espaço contínuo --- atributos em $[0,1]$ ou padronizados, sem normalizar cada
vetor, porque a magnitude carrega informação. Com $\beta\in\{1,10,100,1000\}$ (ou seja, $\sigma=\beta^{-1/2}$),
os escores são $E(\xi^0)$ (Eq.~\ref{eq:energia-moderna}), o resíduo de recuperação
$\|\xi^0-X\softmax(\beta X^{\top}\xi^0)\|$ e a entropia do vetor $\softmax(\beta X^{\top}\xi^0)$
(Algoritmo~\ref{alg:moderna}). O custo por consulta é $O(NK)$ com $N=38$ e $K$ até o número de linhas de treino.
Pela Seção~\ref{sec:moderna}, esperamos que esses escores se comportem como escores de densidade por kernel.

\subsection{Linhas de base}\label{sec:baselines}
O enunciado exige ao menos dois modelos de classificação da lista (\emph{Decision Trees}, \emph{Random Forest},
\emph{Naive Bayes}, SVM, MLP, XGBoost e KNN)~\cite{enunciado}. Como a rede de Hopfield não consta da lista
(Seção~\ref{sec:intro}), tratamo-la como \emph{modelo do tema} e usamos as linhas de base abaixo, implementadas
com o \texttt{scikit-learn}~\cite{pedregosa2011}.
\begin{itemize}
  \item \textbf{Supervisionadas (Etapa~I):} \emph{Random Forest} (com \texttt{class\_weight='balanced'}) e SVM com
        núcleo RBF, treinadas em $L_{\mathrm{tr}}$ (Seção~\ref{sec:divisao}) sobre os 38 atributos originais. São uma
        \emph{referência superior}: usam rótulos, ao contrário do detector \emph{one-class}, e portanto a
        comparação com o Hopfield não é de igual para igual.
  \item \textbf{\emph{One-class} (comparação justa):} \emph{One-Class SVM} e distância ao $k$-ésimo vizinho mais
        próximo (\emph{k}-NN), ajustados apenas na metade de treino sem rótulos, como a rede de Hopfield.
  \item \textbf{Opcional:} MLP (classificador ou autoencoder).
\end{itemize}

\subsection{Divisão dos dados e protocolos de avaliação}\label{sec:divisao}
A metade de treino do SMD não tem rótulos; rótulos --- para treinar modelos supervisionados, escolher
limiares e avaliar --- só existem na metade de teste. Dividimos essa metade \emph{cronologicamente} em três
blocos, nunca por sorteio: pontos vizinhos de uma série temporal são fortemente correlacionados e um sorteio
vazaria informação do teste para o treino. Seja $n$ o número de linhas da metade rotulada,
$c_1=\mathrm{round}(0{,}4\,n)$ e $c_2=\mathrm{round}(0{,}6\,n)$:
\begin{equation}\label{eq:blocos}
L_{\mathrm{tr}}=[0,\,c_1-60),\qquad L_{\mathrm{val}}=[c_1+60,\,c_2-60),\qquad L_{\mathrm{te}}=[c_2+60,\,n),
\end{equation}
com um \emph{embargo} de 60 amostras (1 hora) de cada lado de cada corte, para que nenhuma informação vaze
entre blocos vizinhos. Cada corte é deslocado, se necessário, para a zona mais próxima sem anomalias em
$\pm60$ amostras, de modo que nenhum segmento anômalo atravesse uma fronteira nem caia no embargo. Dos 56
cortes (28 máquinas $\times$ 2), apenas 5 foram deslocados: $c_1$ nas máquinas 2-4 (9.413 em vez de 9.476) e
3-7 (11.431 em vez de 11.482), e $c_2$ nas máquinas 1-1 (16.903 em vez de 17.087), 2-9 (17.346 em vez de 17.233)
e 3-8 (17.336 em vez de 17.222). A metade de treino (sem rótulos) também é dividida em ordem cronológica:
$S_{\mathrm{aj}}$ (primeiros 80\,\%), onde se ajustam os modelos \emph{one-class}, e $S_{\mathrm{cal}}$ (últimos
20\,\%), usada para calibrar o limiar sem rótulos.

\begin{figure}[htbp]
\centering
\begin{tikzpicture}[x=12cm,y=1cm,font=\footnotesize]
  \node[anchor=south west] at (0,3.55) {\textbf{Metade de treino} (sem rótulos; anterior no tempo)};
  \draw[thick,fill=blue!10] (0,2.8) rectangle (0.8,3.5);
  \draw[thick,fill=blue!25] (0.8,2.8) rectangle (1,3.5);
  \node at (0.4,3.15) {$S_{\mathrm{aj}}$ (primeiros 80\,\%)};
  \node at (0.9,3.15) {$S_{\mathrm{cal}}$};
  \node[anchor=north,align=center,text width=4cm,font=\scriptsize] at (0.4,2.75)
    {ajuste dos modelos \emph{one-class}};
  \node[anchor=north,align=center,text width=2.2cm,font=\scriptsize] at (0.9,2.75)
    {calibração do limiar sem rótulos};

  \node[anchor=south west] at (0,1.55) {\textbf{Metade rotulada} (teste do SMD; posterior no tempo)};
  \draw[thick,fill=green!12] (0,0.8) rectangle (0.37,1.5);
  \draw[thick,fill=gray!45] (0.37,0.8) rectangle (0.43,1.5);
  \draw[thick,fill=yellow!30] (0.43,0.8) rectangle (0.57,1.5);
  \draw[thick,fill=gray!45] (0.57,0.8) rectangle (0.63,1.5);
  \draw[thick,fill=red!12] (0.63,0.8) rectangle (1,1.5);
  \node at (0.185,1.15) {$L_{\mathrm{tr}}$};
  \node at (0.5,1.15) {$L_{\mathrm{val}}$};
  \node at (0.815,1.15) {$L_{\mathrm{te}}$};
  \node[anchor=north,align=center,text width=3cm,font=\scriptsize] at (0.185,0.75) {treino dos modelos supervisionados};
  \node[anchor=north,align=center,text width=2.2cm,font=\scriptsize] at (0.5,0.75) {limiar e hiperparâmetros};
  \node[anchor=north,align=center,text width=3cm,font=\scriptsize] at (0.815,0.75) {avaliação final};
\end{tikzpicture}
\caption{Divisão cronológica de uma máquina (esquema fora de escala). As faixas cinza são os embargos de
60 amostras, centrados nos cortes $c_1$ e $c_2$.}
\label{fig:divisao}
\end{figure}

\paragraph{Elegibilidade das máquinas.} Nem toda máquina tem anomalias em todos os blocos
(Tabela~\ref{tab:elegib}; números por máquina no Apêndice~\ref{apx:maquinas}). Para treinar um classificador
supervisionado é preciso ter as duas classes em $L_{\mathrm{tr}}$; para escolher o limiar, anomalias em
$L_{\mathrm{val}}$. Por isso o Protocolo~S, descrito abaixo, usa 18 máquinas (e 23 para os modelos
\emph{one-class}), e as 5 restantes entram apenas no Protocolo~U. Alguns $L_{\mathrm{val}}$ têm pouquíssimas
anomalias (máquina 3-5: 3; 2-2: 11; 3-3: 15; 1-8: 16; 3-4: 20; 1-6: 27), de modo que o limiar não é ajustado
por máquina, e sim globalmente, sobre o $L_{\mathrm{val}}$ agregado.

\begin{table}[htbp]
\centering
\caption{Elegibilidade das máquinas, conforme a existência de anomalias em cada bloco (cálculo próprio).}
\label{tab:elegib}
\small
\begin{tabularx}{\textwidth}{@{}lcLL@{}}
\toprule
Classe & Anomalias em $L_{\mathrm{tr}}$ / $L_{\mathrm{val}}$ / $L_{\mathrm{te}}$ & Máquinas & Uso\\
\midrule
S (18) & $>0$ / $>0$ / $>0$ & 1-3, 1-4, 1-6, 1-7, 1-8, 2-1, 2-3, 2-4, 2-5, 2-6, 2-7, 3-2, 3-3, 3-4, 3-5, 3-6, 3-7, 3-10 & Protocolo~S (supervisionados e \emph{one-class})\\
O (5) & $0$ / $>0$ / $>0$ & 1-1, 1-5, 2-2, 2-9, 3-8 & Protocolo~S, só \emph{one-class}\\
-- (5) & $L_{\mathrm{val}}=0$ & 1-2, 2-8, 3-1, 3-9, 3-11 & só Protocolo~U\\
\bottomrule
\end{tabularx}
\end{table}

\begin{table}[htbp]
\centering
\caption{Blocos cronológicos agregados (soma das máquinas de cada conjunto; cálculo próprio).}
\label{tab:blocos}
\small
\begin{tabular}{llrrr}
\toprule
Conjunto & Bloco & Linhas & Anômalas & Prevalência\\
\midrule
\multirow{3}{*}{18 máquinas (S)} & $L_{\mathrm{tr}}$  & 175.448 & 6.664  & 3,80\,\%\\
                                 & $L_{\mathrm{val}}$ & 86.273  & 1.970  & 2,28\,\%\\
                                 & $L_{\mathrm{te}}$  & 175.562 & 10.545 & 6,01\,\%\\
\midrule
\multirow{3}{*}{23 máquinas (S+O)} & $L_{\mathrm{tr}}$  & 228.473 & 6.664  & 2,92\,\%\\
                                   & $L_{\mathrm{val}}$ & 112.377 & 4.500  & 4,00\,\%\\
                                   & $L_{\mathrm{te}}$  & 228.544 & 16.768 & 7,34\,\%\\
\bottomrule
\end{tabular}
\end{table}

A prevalência muda de bloco para bloco (de 2,28\,\% a 6,01\,\% nas 18 máquinas), de modo que um limiar ótimo em
$L_{\mathrm{val}}$ não é necessariamente ótimo em $L_{\mathrm{te}}$; esse é um dos motivos para reportar também
métricas independentes de limiar (AUC-ROC e PR-AUC). Para análises detalhadas e para a Etapa~II escolhemos uma
máquina representativa por grupo, todas com anomalias em todos os blocos: \textbf{1-3}, \textbf{2-5} e \textbf{3-6}
(anomalias em $L_{\mathrm{tr}}$/$L_{\mathrm{val}}$/$L_{\mathrm{te}}$: 86/281/450, 314/39/627 e 174/437/583;
prevalências de 3,45\,\%, 4,14\,\% e 4,16\,\%, próximas dos 4,16\,\% do conjunto).

\paragraph{Limiar de decisão.} Escores de máquinas diferentes não são diretamente comparáveis. Para os modelos
\emph{one-class}, o limiar de cada máquina é o quantil $(1-\alpha)$ dos escores em $S_{\mathrm{cal}}$ ---
isto é, $\alpha$ é a taxa de alarmes tolerada sobre o comportamento normal --- e $\alpha$ é um parâmetro
\emph{global}. Para os modelos supervisionados, o parâmetro global é o limiar sobre a probabilidade prevista.

\begin{table}[htbp]
\centering
\caption{Protocolos de avaliação (propostas de projeto).}
\label{tab:protocolos}
\small
\begin{tabularx}{\textwidth}{@{}lLL@{}}
\toprule
 & Protocolo~S (principal) & Protocolo~U (secundário, sem rótulos)\\
\midrule
Máquinas & 18 (classe S); modelos \emph{one-class} também nas 5 da classe O & 28\\
Ajuste & \emph{one-class}: $S_{\mathrm{aj}}$; supervisionados: $L_{\mathrm{tr}}$ & \emph{one-class}: $S_{\mathrm{aj}}$\\
Limiar e hiperparâmetros & globais, escolhidos por $F_1$ em $L_{\mathrm{val}}$ agregado ($\alpha$ ou limiar de probabilidade) & $\alpha\in\{0{,}5;\,1;\,2;\,5\}\,\%$ fixado \emph{a priori}; limiar por máquina em $S_{\mathrm{cal}}$\\
Avaliação & $L_{\mathrm{te}}$ & toda a metade rotulada\\
Relato & todas as métricas da Seção~\ref{sec:metricas} & AUC-ROC, PR-AUC e métricas no limiar sem rótulos; $F_1$ ótimo (``oráculo'') só como teto otimista\\
\bottomrule
\end{tabularx}
\end{table}

\paragraph{Agregação e significância.} Reportamos a média entre máquinas (\emph{macro}) e as contagens
agregadas (\emph{micro}); a AUC-ROC e a PR-AUC são calculadas por máquina e promediadas. Comparações entre
modelos usam o teste de Wilcoxon dos postos sinalizados, pareado por máquina, e as partes estocásticas (ordem de
atualização, inicialização do \emph{k-medoids}, \emph{Random Forest}, MLP) são repetidas com 3 sementes.

\subsection{Balanceamento de classes}\label{sec:balanceamento}
Com 4,16\,\% de pontos anômalos, o desbalanceamento é forte, mas seu tratamento depende do modelo:
\begin{itemize}
  \item o detector \emph{one-class} (Hopfield, OC-SVM, \emph{k}-NN) não usa rótulos no ajuste, logo não há o que
        balancear;
  \item nos modelos supervisionados usamos pesos de classe (\texttt{class\_weight='balanced'}) ou
        subamostragem da classe majoritária, \emph{apenas em $L_{\mathrm{tr}}$};
  \item não usamos SMOTE nem sobreamostragem sintética: elas criam observações que nunca ocorreram e ignoram a
        dependência temporal;
  \item $L_{\mathrm{val}}$ e $L_{\mathrm{te}}$ mantêm a prevalência natural e nunca são balanceados (a única
        exceção, a subamostra da Etapa~II, é corrigida por prevalência; Seção~\ref{sec:metricas}).
\end{itemize}

\subsection{Métricas de avaliação}\label{sec:metricas}
Com a classe anômala como positiva e $TP$, $FP$, $FN$, $TN$ as contagens ponto a ponto,
\begin{align*}
\text{acurácia}&=\frac{TP+TN}{TP+TN+FP+FN}, & P&=\frac{TP}{TP+FP}, & R&=\frac{TP}{TP+FN},\\
F_\gamma&=(1+\gamma^2)\frac{P\,R}{\gamma^2P+R}, & \mathrm{FPR}&=\frac{FP}{FP+TN},
\end{align*}
\[
\mathrm{MCC}=\frac{TP\cdot TN-FP\cdot FN}{\sqrt{(TP+FP)(TP+FN)(TN+FP)(TN+FN)}} .
\]
Reportamos as métricas mínimas do enunciado: acurácia, precisão ($P$), revocação ($R$, \emph{recall}), $F_1$ e
$F_2$ ($\gamma=1$ e $\gamma=2$; este pesa mais a revocação), AUC-ROC~\cite{fawcett2006} e matriz de confusão.
Como complemento, a PR-AUC (preferível à AUC-ROC com classes desbalanceadas~\cite{saito2015}), o MCC, o tempo
de inferência por 100 amostras e o custo de memória ($O(N^2)$ na rede clássica). O limiar é escolhido
maximizando $F_1$ em $L_{\mathrm{val}}$ agregado; $F_2$ é reportado no mesmo limiar.

\paragraph{Avaliação por ponto e por segmento.} A avaliação principal é \emph{por ponto}. Na avaliação por
segmento (\emph{point-adjust}), basta acertar um ponto de um segmento anômalo para que o segmento inteiro conte como
acerto, o que infla as métricas~\cite{kim2022}; ela é a usada em parte da literatura do SMD~\cite{su2019} e a
reportamos apenas como secundária, para comparação.

\paragraph{Precisão ajustada à prevalência.} Se um conjunto for subamostrado com prevalência diferente da natural
$\pi$ (caso da Etapa~II), a precisão passa a ser $P_{\mathrm{nat}}=\dfrac{R\,\pi}{R\,\pi+\mathrm{FPR}\,(1-\pi)}$,
pois $R$ e $\mathrm{FPR}$ não dependem da prevalência.

\subsection{Experimentos e hipóteses}\label{sec:experimentos}
\begin{description}
  \item[E1a --- capacidade em padrões sintéticos.] Para $N\in\{67,122,224,412\}$, armazenar $K$ padrões
        aleatórios não viesados (Hebb e Storkey) e medir a fração de padrões que são pontos fixos e a taxa de
        recuperação a partir de versões ruidosas (inversão de uma fração dos bits) em função de $K/N$,
        comparando com a Tabela~\ref{tab:capacidade}. Valida a implementação.
  \item[E1b --- capacidade com protótipos do SMD.] A mesma análise com os protótipos codificados do SMD,
        para ver o efeito de bits viesados e correlacionados.
  \item[E2 --- detecção de anomalias.] Protocolos~S e~U (Seção~\ref{sec:divisao}), com os modelos da
        Seção~\ref{sec:baselines}.
  \item[E3 --- Etapa~II.] Protocolo da Seção~\ref{sec:etapa2}.
  \item[E4 (opcional) --- interpretabilidade.] Comparar os atributos cujos bits foram alterados na recuperação
        com os atributos indicados em \texttt{interpretation\_label}.
\end{description}
\noindent As \textbf{ablações} cobrem: número de níveis $b\in\{2,4,8,16\}$; número de protótipos
$K\in\{2,5,10,20,30\}$; regra (Hebb ou Storkey); dinâmica recorrente ($S_1$ \emph{versus} $S_2$); janelas
temporais; $\beta\in\{1,10,100,1000\}$ na variante moderna; e, de forma exploratória e fora da bibliografia
fornecida, uma regra de Hebb \emph{centrada} para bits viesados.

\paragraph{Hipóteses (expectativas \emph{a priori}, não resultados).}
\begin{description}
  \item[H1] Com a regra de Hebb, a recuperação se degrada quando $K/N$ passa de cerca de $0{,}1$.
  \item[H2] A regra de Storkey é ao menos tão boa quanto a de Hebb com protótipos correlacionados.
  \item[H3] A rede moderna, com $\beta$ adequado, detecta melhor que a clássica.
  \item[H4] No Protocolo~S, \emph{Random Forest} e SVM superam o Hopfield \emph{one-class}; no Protocolo~U, o
        Hopfield é competitivo com OC-SVM e \emph{k}-NN.
  \item[H5] O LLM em \emph{zero-shot} é fraco em telemetria numérica anônima; o \emph{few-shot} melhora, mas não
        supera o \emph{Random Forest}.
\end{description}

\subsection{Etapa II: LLM \emph{versus} modelos clássicos}\label{sec:etapa2}
A Etapa~II compara um modelo de linguagem usado por \emph{prompting} com os modelos clássicos da Etapa~I,
sobre \emph{os mesmos dados}.
\begin{itemize}
  \item \textbf{Tarefa:} dada uma observação (as 38 métricas anônimas em $[0,1]$, serializadas como texto),
        responder em JSON o rótulo (0 ou 1) e um escore de anomalia em $[0,1]$. Temperatura 0. Registrar o
        modelo, a versão e a data de execução.
  \item \textbf{Configurações:} \emph{zero-shot} e \emph{few-shot} com $k\in\{3,5,10,20\}$ exemplos de
        $L_{\mathrm{tr}}$, aproximadamente balanceados (2N+1A, 3N+2A, 5N+5A, 10N+10A, em que N = normal e
        A = anômalo), sorteados com 3 sementes.
  \item \textbf{Dados:} uma subamostra de $L_{\mathrm{te}}$ (500 a 2.000 pontos, conforme o custo) nas máquinas
        representativas 1-3, 2-5 e 3-6. Os modelos da Etapa~I são reavaliados nessa mesma subamostra, com a
        precisão ajustada por prevalência.
  \item \textbf{Métricas:} as da Etapa~I (Seção~\ref{sec:metricas}), mais o tempo por 100 amostras e o custo
        estimado de API por 100 amostras (preço por \emph{token} $\times$ \emph{tokens} do \emph{prompt} e da
        resposta; preços a preencher).
  \item \textbf{\emph{Prompt}:} Apêndice~\ref{apx:prompt}. Variante opcional: \emph{zero-shot} com um perfil de
        referência (mediana e percentil 95 de cada métrica no treino normal).
\end{itemize}
\noindent Limitações esperadas: as métricas são anônimas, de modo que o LLM não tem conhecimento prévio sobre o
significado de cada uma; o escore autodeclarado do LLM não é calibrado; o resultado é sensível ao formato e à
ordem dos exemplos; há custo e latência; e a reprodutibilidade é limitada pela atualização dos modelos. Esta
versão do relatório define o protocolo; a execução depende de acesso a um modelo de linguagem
(Seção~\ref{sec:resultados}).

\subsection{Reprodutibilidade}\label{sec:repro}
\begin{itemize}
  \item \textbf{Dados:} obtidos do repositório do OmniAnomaly~\cite{omnianomaly}, pasta
        \texttt{ServerMachineDataset}, com os subdiretórios \texttt{train}, \texttt{test},
        \texttt{test\_label} e \texttt{interpretation\_label} e um arquivo \texttt{machine-G-I.txt} por
        máquina (CSV sem cabeçalho, 38 colunas; um rótulo por linha). Registrar o \emph{commit} usado.
  \item \textbf{Ambiente:} Python 3, NumPy e \texttt{scikit-learn}~\cite{pedregosa2011}; registrar as versões.
  \item \textbf{Sementes:} 3 por experimento estocástico, fixadas e registradas.
  \item \textbf{Registros:} salvar os escores por máquina e por ponto, para recalcular as métricas sem reexecutar
        os modelos; medir os tempos sempre na mesma máquina.
  \item \textbf{Código:} os \emph{scripts} dos experimentos não fazem parte desta versão do relatório.
\end{itemize}

\section{Resultados}\label{sec:resultados}

\begin{notabox}{Estado desta seção}
Esta versão do relatório \textbf{não contém resultados de modelos}. As tabelas abaixo fixam o formato de relato;
as células marcadas com \lacuna\ serão preenchidas após os experimentos --- nenhum valor foi estimado ou
inventado. As estatísticas já obtidas, que são reais, estão nas Tabelas~\ref{tab:grupos}, \ref{tab:exemplos},
\ref{tab:deriva}, \ref{tab:codificacao}, \ref{tab:elegib} e~\ref{tab:blocos} e no Apêndice~\ref{apx:maquinas}.
\end{notabox}

\subsection{E1: capacidade e recuperação}\label{sec:res-e1}
A Tabela~\ref{tab:res-e1a} reúne a validação da implementação em padrões sintéticos (E1a). O experimento E1b
usa o mesmo formato com os protótipos do SMD, para $K\in\{2,5,10,20,30\}$ e $b\in\{2,4,8,16\}$. Esperamos que
a fração de pontos fixos caia perto de $K/N\approx0{,}14$ para a regra de Hebb (Tabela~\ref{tab:capacidade}).

\begin{table}[htbp]
\centering
\caption{E1a --- fração de padrões aleatórios, não viesados, que são pontos fixos (\%), por $N$, razão $K/N$ e
regra de armazenamento (H = Hebb; S = Storkey). A preencher.}
\label{tab:res-e1a}
\footnotesize
\setlength{\tabcolsep}{4pt}
\begin{tabular}{rcccccccc}
\toprule
 & \multicolumn{2}{c}{$N=67$} & \multicolumn{2}{c}{$N=122$} & \multicolumn{2}{c}{$N=224$} & \multicolumn{2}{c}{$N=412$}\\
\cmidrule(lr){2-3}\cmidrule(lr){4-5}\cmidrule(lr){6-7}\cmidrule(lr){8-9}
$K/N$ & H & S & H & S & H & S & H & S\\
\midrule
0,05  & \lacuna & \lacuna & \lacuna & \lacuna & \lacuna & \lacuna & \lacuna & \lacuna\\
0,10  & \lacuna & \lacuna & \lacuna & \lacuna & \lacuna & \lacuna & \lacuna & \lacuna\\
0,138 & \lacuna & \lacuna & \lacuna & \lacuna & \lacuna & \lacuna & \lacuna & \lacuna\\
0,20  & \lacuna & \lacuna & \lacuna & \lacuna & \lacuna & \lacuna & \lacuna & \lacuna\\
0,30  & \lacuna & \lacuna & \lacuna & \lacuna & \lacuna & \lacuna & \lacuna & \lacuna\\
\bottomrule
\end{tabular}
\end{table}

\subsection{E2: detecção de anomalias}\label{sec:res-e2}
\paragraph{Protocolo S.} A Tabela~\ref{tab:res-s} apresenta as métricas médias entre as 18 máquinas da
classe S em $L_{\mathrm{te}}$ (macro; a AUC-ROC e a PR-AUC são calculadas por máquina e promediadas), e a
Tabela~\ref{tab:res-conf} as matrizes de confusão agregadas (micro). Os modelos supervisionados são uma
referência superior (Seção~\ref{sec:baselines}).

\begin{table}[htbp]
\centering
\caption{Protocolo S --- métricas médias entre máquinas em $L_{\mathrm{te}}$ (macro). Limiar escolhido por $F_1$
em $L_{\mathrm{val}}$ agregado. Tempo: inferência por 100 amostras. A preencher.}
\label{tab:res-s}
\footnotesize
\setlength{\tabcolsep}{3pt}
\begin{tabularx}{\textwidth}{@{}Lccccccccc@{}}
\toprule
Modelo & Acur. & $P$ & $R$ & $F_1$ & $F_2$ & AUC-ROC & PR-AUC & MCC & s/100\\
\midrule
Hopfield clássico (Hebb)   & \lacuna & \lacuna & \lacuna & \lacuna & \lacuna & \lacuna & \lacuna & \lacuna & \lacuna\\
Hopfield clássico (Storkey)& \lacuna & \lacuna & \lacuna & \lacuna & \lacuna & \lacuna & \lacuna & \lacuna & \lacuna\\
Hopfield moderno           & \lacuna & \lacuna & \lacuna & \lacuna & \lacuna & \lacuna & \lacuna & \lacuna & \lacuna\\
OC-SVM                     & \lacuna & \lacuna & \lacuna & \lacuna & \lacuna & \lacuna & \lacuna & \lacuna & \lacuna\\
\emph{k}-NN (dist. ao $k$-ésimo vizinho) & \lacuna & \lacuna & \lacuna & \lacuna & \lacuna & \lacuna & \lacuna & \lacuna & \lacuna\\
\emph{Random Forest} (supervisionado) & \lacuna & \lacuna & \lacuna & \lacuna & \lacuna & \lacuna & \lacuna & \lacuna & \lacuna\\
SVM-RBF (supervisionado)   & \lacuna & \lacuna & \lacuna & \lacuna & \lacuna & \lacuna & \lacuna & \lacuna & \lacuna\\
\bottomrule
\end{tabularx}
\end{table}

\begin{table}[htbp]
\centering
\caption{Protocolo S --- matrizes de confusão agregadas (micro) em $L_{\mathrm{te}}$ (positivo = anômalo). A
preencher.}
\label{tab:res-conf}
\small
\begin{tabular}{lrrrr}
\toprule
Modelo & $TP$ & $FP$ & $FN$ & $TN$\\
\midrule
Hopfield clássico (melhor configuração) & \lacuna & \lacuna & \lacuna & \lacuna\\
Hopfield moderno                        & \lacuna & \lacuna & \lacuna & \lacuna\\
OC-SVM                                  & \lacuna & \lacuna & \lacuna & \lacuna\\
\emph{k}-NN                             & \lacuna & \lacuna & \lacuna & \lacuna\\
\emph{Random Forest}                    & \lacuna & \lacuna & \lacuna & \lacuna\\
SVM-RBF                                 & \lacuna & \lacuna & \lacuna & \lacuna\\
\bottomrule
\end{tabular}
\end{table}

\paragraph{Protocolo U.} A Tabela~\ref{tab:res-u} traz os modelos \emph{one-class} avaliados nas 28 máquinas,
com limiar sem rótulos definido por $\alpha$. A AUC-ROC e a PR-AUC independem de $\alpha$.

\begin{table}[htbp]
\centering
\caption{Protocolo U --- 28 máquinas, toda a metade rotulada, limiar sem rótulos (quantil $1-\alpha$ em
$S_{\mathrm{cal}}$). $F_1^{\ast}$: $F_1$ ótimo com acesso aos rótulos (teto otimista, não comparável). A
preencher.}
\label{tab:res-u}
\footnotesize
\setlength{\tabcolsep}{4pt}
\begin{tabular}{llccccccc}
\toprule
Modelo & $\alpha$ & Acur. & $P$ & $R$ & $F_1$ & AUC-ROC & PR-AUC & $F_1^{\ast}$\\
\midrule
\multirow{4}{*}{Hopfield clássico} & 0,5\,\% & \lacuna & \lacuna & \lacuna & \lacuna & \multirow{4}{*}{\lacuna} & \multirow{4}{*}{\lacuna} & \multirow{4}{*}{\lacuna}\\
 & 1\,\% & \lacuna & \lacuna & \lacuna & \lacuna & & & \\
 & 2\,\% & \lacuna & \lacuna & \lacuna & \lacuna & & & \\
 & 5\,\% & \lacuna & \lacuna & \lacuna & \lacuna & & & \\
\midrule
\multirow{4}{*}{Hopfield moderno} & 0,5\,\% & \lacuna & \lacuna & \lacuna & \lacuna & \multirow{4}{*}{\lacuna} & \multirow{4}{*}{\lacuna} & \multirow{4}{*}{\lacuna}\\
 & 1\,\% & \lacuna & \lacuna & \lacuna & \lacuna & & & \\
 & 2\,\% & \lacuna & \lacuna & \lacuna & \lacuna & & & \\
 & 5\,\% & \lacuna & \lacuna & \lacuna & \lacuna & & & \\
\midrule
\multirow{4}{*}{OC-SVM} & 0,5\,\% & \lacuna & \lacuna & \lacuna & \lacuna & \multirow{4}{*}{\lacuna} & \multirow{4}{*}{\lacuna} & \multirow{4}{*}{\lacuna}\\
 & 1\,\% & \lacuna & \lacuna & \lacuna & \lacuna & & & \\
 & 2\,\% & \lacuna & \lacuna & \lacuna & \lacuna & & & \\
 & 5\,\% & \lacuna & \lacuna & \lacuna & \lacuna & & & \\
\midrule
\multirow{4}{*}{\emph{k}-NN} & 0,5\,\% & \lacuna & \lacuna & \lacuna & \lacuna & \multirow{4}{*}{\lacuna} & \multirow{4}{*}{\lacuna} & \multirow{4}{*}{\lacuna}\\
 & 1\,\% & \lacuna & \lacuna & \lacuna & \lacuna & & & \\
 & 2\,\% & \lacuna & \lacuna & \lacuna & \lacuna & & & \\
 & 5\,\% & \lacuna & \lacuna & \lacuna & \lacuna & & & \\
\bottomrule
\end{tabular}
\end{table}

\subsection{Ablações do detector de Hopfield}\label{sec:res-abl}
A Tabela~\ref{tab:res-abl} mostra o efeito do número de níveis $b$, do número de protótipos $K$ e da regra de
armazenamento sobre a PR-AUC média (Protocolo~S, escore $S_1$), e a Tabela~\ref{tab:res-escores} compara os
escores da Tabela~\ref{tab:escores}.

\begin{table}[htbp]
\centering
\caption{Ablação --- PR-AUC média (macro) do detector clássico, escore $S_1$, Protocolo~S, por $b$, $K$ e regra.
A preencher.}
\label{tab:res-abl}
\small
\setlength{\tabcolsep}{5pt}
\begin{tabular}{llccccc}
\toprule
Regra & $b$ & $K=2$ & $K=5$ & $K=10$ & $K=20$ & $K=30$\\
\midrule
\multirow{4}{*}{Hebb}    & 2  & \lacuna & \lacuna & \lacuna & \lacuna & \lacuna\\
                         & 4  & \lacuna & \lacuna & \lacuna & \lacuna & \lacuna\\
                         & 8  & \lacuna & \lacuna & \lacuna & \lacuna & \lacuna\\
                         & 16 & \lacuna & \lacuna & \lacuna & \lacuna & \lacuna\\
\midrule
\multirow{4}{*}{Storkey} & 2  & \lacuna & \lacuna & \lacuna & \lacuna & \lacuna\\
                         & 4  & \lacuna & \lacuna & \lacuna & \lacuna & \lacuna\\
                         & 8  & \lacuna & \lacuna & \lacuna & \lacuna & \lacuna\\
                         & 16 & \lacuna & \lacuna & \lacuna & \lacuna & \lacuna\\
\bottomrule
\end{tabular}
\end{table}

\begin{table}[htbp]
\centering
\caption{Comparação dos escores $S_1$--$S_6$ na melhor configuração do detector clássico (Protocolo~S, macro).
A preencher.}
\label{tab:res-escores}
\small
\begin{tabular}{lcc}
\toprule
Escore & AUC-ROC & PR-AUC\\
\midrule
$S_1$ & \lacuna & \lacuna\\
$S_2$ & \lacuna & \lacuna\\
$S_3$ & \lacuna & \lacuna\\
$S_4$ & \lacuna & \lacuna\\
$S_5$ & \lacuna & \lacuna\\
$S_6$ & \lacuna & \lacuna\\
\bottomrule
\end{tabular}
\end{table}

\subsection{E3: Etapa II --- LLM \emph{versus} modelos clássicos}\label{sec:res-e3}
A Tabela~\ref{tab:res-llm} segue o protocolo da Seção~\ref{sec:etapa2}. Os modelos clássicos são reavaliados na
\emph{mesma} subamostra.

\begin{table}[htbp]
\centering
\caption{Etapa II --- LLM (modelo, versão e data: \apreencher) e modelos clássicos na mesma subamostra de
$L_{\mathrm{te}}$ (máquinas 1-3, 2-5 e 3-6; $n=\ $\apreencher). Precisão ajustada por prevalência. Custo:
estimativa de API por 100 amostras (preços \apreencher). A preencher.}
\label{tab:res-llm}
\footnotesize
\setlength{\tabcolsep}{3pt}
\begin{tabularx}{\textwidth}{@{}Lcccccccc@{}}
\toprule
Configuração & Acur. & $P_{\mathrm{nat}}$ & $R$ & $F_1$ & $F_2$ & AUC-ROC & s/100 & Custo/100\\
\midrule
LLM, \emph{zero-shot}       & \lacuna & \lacuna & \lacuna & \lacuna & \lacuna & \lacuna & \lacuna & \lacuna\\
LLM, \emph{few-shot} $k=3$  & \lacuna & \lacuna & \lacuna & \lacuna & \lacuna & \lacuna & \lacuna & \lacuna\\
LLM, \emph{few-shot} $k=5$  & \lacuna & \lacuna & \lacuna & \lacuna & \lacuna & \lacuna & \lacuna & \lacuna\\
LLM, \emph{few-shot} $k=10$ & \lacuna & \lacuna & \lacuna & \lacuna & \lacuna & \lacuna & \lacuna & \lacuna\\
LLM, \emph{few-shot} $k=20$ & \lacuna & \lacuna & \lacuna & \lacuna & \lacuna & \lacuna & \lacuna & \lacuna\\
\midrule
\emph{Random Forest}        & \lacuna & \lacuna & \lacuna & \lacuna & \lacuna & \lacuna & \lacuna & ---\\
SVM-RBF                     & \lacuna & \lacuna & \lacuna & \lacuna & \lacuna & \lacuna & \lacuna & ---\\
Hopfield clássico           & \lacuna & \lacuna & \lacuna & \lacuna & \lacuna & \lacuna & \lacuna & ---\\
OC-SVM                      & \lacuna & \lacuna & \lacuna & \lacuna & \lacuna & \lacuna & \lacuna & ---\\
\bottomrule
\end{tabularx}
\end{table}

\subsection{E4 (opcional): interpretabilidade}
Quando realizado, comparará os atributos cujos bits foram alterados na recuperação com os atributos de
\texttt{interpretation\_label}: \apreencher.

\section{Comparação crítica}\label{sec:critica}

Esta seção organiza a análise crítica em dois eixos: as hipóteses formuladas \emph{antes} dos experimentos
(Seção~\ref{sec:experimentos}), com a evidência necessária para aceitá-las ou rejeitá-las, e os pontos
positivos e negativos da abordagem, separando o que decorre da teoria e dos dados já examinados do que só os
experimentos podem mostrar.

\subsection{Hipóteses e evidência necessária}
\begin{table}[htbp]
\centering
\caption{Hipóteses, evidência necessária e resultado empírico (a preencher após os experimentos).}
\label{tab:hipoteses}
\small
\begin{tabularx}{\textwidth}{@{}c>{\hsize=1.15\hsize}L>{\hsize=.85\hsize}Lcc@{}}
\toprule
 & Hipótese (resumo) & Evidência necessária & Resultado & Veredito\\
\midrule
H1 & Hebb degrada com $K/N\gtrsim0{,}1$ & E1a e E1b (Tabela~\ref{tab:res-e1a}) & \lacuna & \lacuna\\
H2 & Storkey $\ge$ Hebb com protótipos correlacionados & E1a, E1b e ablação (Tabela~\ref{tab:res-abl}) & \lacuna & \lacuna\\
H3 & Moderna $>$ clássica & Protocolos~S e~U (Tabelas~\ref{tab:res-s} e~\ref{tab:res-u}); teste de Wilcoxon & \lacuna & \lacuna\\
H4 & S: supervisionados $>$ Hopfield; U: Hopfield $\approx$ OC-SVM e \emph{k}-NN & Tabelas~\ref{tab:res-s} e~\ref{tab:res-u}; teste de Wilcoxon & \lacuna & \lacuna\\
H5 & LLM \emph{zero-shot} fraco; \emph{few-shot} melhora, mas não supera o \emph{Random Forest} & Etapa~II (Tabela~\ref{tab:res-llm}) & \lacuna & \lacuna\\
\bottomrule
\end{tabularx}
\end{table}

\paragraph{Critérios definidos antes dos experimentos (proposta).}
\begin{itemize}
  \item A PR-AUC de um detector aleatório é aproximadamente a prevalência (cerca de 4,16\,\% no conjunto);
        um detector só é considerado informativo se superar esse patamar de forma consistente entre as máquinas.
  \item O detector de Hopfield é considerado \emph{competitivo} se, no Protocolo~U, sua PR-AUC média não for
        significativamente inferior à do melhor entre OC-SVM e \emph{k}-NN (Wilcoxon pareado por máquina, nível
        de 5\,\%).
  \item Diferenças entre configurações só são destacadas quando maiores que a variação entre as 3 sementes.
\end{itemize}

\subsection{Pontos positivos e negativos}
A Tabela~\ref{tab:prospros} lista, por aspecto, o que a teoria e os dados já examinados permitem esperar; a
última coluna é reservada às evidências dos experimentos.

{\small
\begin{longtable}{@{}P{2.8cm}P{4.3cm}P{4.3cm}P{3.1cm}@{}}
\caption{Pontos positivos e negativos esperados (\emph{a priori}) e evidência empírica (a preencher).}
\label{tab:prospros}\\
\toprule
Aspecto & Positivo (\emph{a priori}) & Negativo (\emph{a priori}) & Evidência empírica\\
\midrule
\endfirsthead
\toprule
Aspecto & Positivo (\emph{a priori}) & Negativo (\emph{a priori}) & Evidência empírica\\
\midrule
\endhead
\midrule
\multicolumn{4}{r}{\emph{continua na próxima página}}\\
\endfoot
\bottomrule
\endlastfoot
Treinamento &
Regras locais e incrementais (Hebb, Storkey): um novo padrão é somado sem revisitar os anteriores~\cite{wikihopfield}. Treino \emph{one-class}, sem rótulos, compatível com a metade de treino do SMD. &
O treino não rotulado pode conter anomalias, que viram protótipos ``normais'' (contaminação). &
\lacuna\\[2pt]
Capacidade &
Para padrões aleatórios, a rede moderna tem capacidade exponencial~\cite{ramsauer2021}. &
A rede clássica guarda ${\approx}0{,}138N$ padrões~\cite{hertz1991}; bits do termômetro viesados reduzem a capacidade efetiva (Tabela~\ref{tab:codificacao}); poucos protótipos podem não cobrir um comportamento normal multimodal. &
\lacuna\\[2pt]
Estados espúrios &
O escore $S_2$ (distância ao protótipo mais próximo) detecta o padrão invertido $-\xi$ (Apêndice~\ref{apx:exemplo}). &
Mínimos espúrios, em número exponencial~\cite{bruckroychowdhury1990,ramsauer2021}, podem capturar uma observação anômala e mascará-la. &
\lacuna\\[2pt]
Representação &
A codificação em termômetro preserva a ordem dos valores e dá uma distância de Hamming ordinal (Apêndice~\ref{apx:deriv}). &
A binarização perde informação e multiplica a dimensão: $N$ de 67 a 412 para 38 atributos. &
\lacuna\\[2pt]
Tempo e deriva &
Inferência rápida e com memória fixa: $O(N^2)$ valores na rede clássica. &
Sem modelagem temporal (a menos de janelas); sensível à deriva entre treino e teste, que é heterogênea (Tabela~\ref{tab:deriva}); exige recalibração do limiar. &
\lacuna\\[2pt]
Custo computacional &
Rede clássica: poucas operações por varredura. &
A rede moderna guarda os $K$ padrões, com custo $O(NK)$ por consulta e $O(NK)$ de memória. &
\lacuna\\[2pt]
Interpreta\-bilidade &
Os bits alterados na recuperação apontam atributos suspeitos (E4). &
Atributos anônimos: sem significado para o operador. &
\lacuna\\[2pt]
Rede moderna &
Recuperação em uma atualização; ligação com a atenção do \emph{Transformer}~\cite{ramsauer2021}. &
A energia equivale a um escore de densidade por kernel (Seção~\ref{sec:moderna}): pouca vantagem esperada sobre \emph{k}-NN e KDE. &
\lacuna\\[2pt]
Adoção &
Base teórica clara (energia, convergência). &
A rede clássica é raramente usada em aprendizado de máquina moderno~\cite{wikihopfield}; não consta da lista de modelos do enunciado. &
\lacuna\\
\end{longtable}
}

\subsection{Síntese a priori}
Se as hipóteses se confirmarem, o quadro esperado é: o Hopfield \emph{one-class} é útil onde rótulos são
escassos ou inexistentes e a latência importa, mas não deve superar um classificador supervisionado quando há
rótulos (H4); a variante moderna deve dar ganho sobre a clássica, porém semelhante ao de métodos de densidade
(H3); e o LLM sem conhecimento prévio sobre métricas anônimas deve ficar abaixo dos modelos treinados (H5). Essas
expectativas são \emph{hipóteses}: a Seção~\ref{sec:conclusoes} indica como serão revisadas diante dos
resultados.

\section{Conclusões}\label{sec:conclusoes}

\subsection{O que já pode ser afirmado}
O texto desta versão sustenta as afirmações abaixo, baseadas na teoria e na exploração dos dados; nenhuma delas
depende de resultados de modelos.
\begin{itemize}
  \item A rede de Hopfield discreta é uma memória associativa cuja dinâmica assíncrona não aumenta a energia e
        termina em um ponto fixo (Proposição~\ref{prop:conv}); isso a torna um filtro de recuperação bem
        definido, mas com capacidade limitada (${\approx}0{,}138N$) e estados espúrios, entre eles o padrão
        invertido $-\xi$ (Seção~\ref{sec:capacidade}).
  \item Formulamos um detector \emph{one-class} sobre essa rede, com escores baseados em distância de Hamming e
        em energia (Seção~\ref{sec:escores}). Mostramos, por derivação própria, que o escore de energia da rede
        moderna é um mínimo suave de distâncias, equivalente a uma densidade por kernel
        (Seção~\ref{sec:moderna}); por isso não esperamos vantagem automática sobre \emph{k}-NN ou KDE.
  \item No SMD, apenas 4,16\,\% dos pontos de teste são anômalos; um classificador que sempre responde
        ``normal'' acerta 95,84\,\%, o que torna a acurácia, isoladamente, uma métrica inadequada
        (Seção~\ref{sec:smd}).
  \item Os dados têm atributos constantes no treino que passam a variar em anomalias, deriva heterogênea entre
        treino e teste e muitas anomalias sutis, dentro da faixa do treino; ``acima do máximo'' sozinho é um
        detector pobre (Seção~\ref{sec:exploracao}).
  \item A codificação em termômetro produz de 67 a 412 bits e bits viesados; armazenar todos os pontos de
        treino é inviável (de 501 a 3.080 vezes a capacidade clássica), o que exige poucos protótipos
        (Seção~\ref{sec:codificacao}).
  \item Só 18 das 28 máquinas permitem um protocolo supervisionado completo (23 para modelos \emph{one-class});
        o Protocolo~U cobre as 28 (Seção~\ref{sec:divisao}).
\end{itemize}

\subsection{O que depende dos experimentos}
Dependem de execução: se o detector de Hopfield separa pontos normais de anômalos acima do acaso e de métodos
\emph{one-class} simples; o efeito de $b$, $K$ e da regra de armazenamento; se a dinâmica recorrente
acrescenta algo à busca do protótipo mais próximo; se a rede moderna supera a clássica; e o comportamento do LLM
na Etapa~II. As hipóteses H1--H5 e os critérios de decisão estão na Seção~\ref{sec:critica}; os resultados
serão registrados na Seção~\ref{sec:resultados}.

\subsection{Limitações}
\begin{itemize}
  \item \textbf{Representação:} a binarização perde informação; bits correlacionados e viesados reduzem a
        capacidade efetiva; poucos protótipos podem não cobrir um comportamento normal multimodal.
  \item \textbf{Modelo:} estados espúrios podem mascarar anomalias; a memória da rede clássica cresce com
        $N^2$; não há modelagem temporal, a menos de janelas; a rede é sensível a hiperparâmetros e ao limiar.
  \item \textbf{Dados:} atributos anônimos e de ``baixo nível'' apenas aproximado; rótulos de incidentes com
        fronteiras imprecisas; treino sem rótulos e possivelmente contaminado; deriva entre treino e teste;
        grande variação entre máquinas.
  \item \textbf{Escopo:} esta versão não traz resultados experimentais nem a execução da Etapa~II, e o Hopfield
        não consta da lista de modelos do enunciado (a confirmar com o professor).
\end{itemize}

\subsection{Próximos passos}
\begin{enumerate}
  \item Confirmar com o professor o uso da rede de Hopfield como modelo do tema, com \emph{Random Forest} e SVM
        como linhas de base da Etapa~I.
  \item Implementar o \emph{pipeline} (codificação, protótipos, Hebb e Storkey, recuperação, escores) e
        validá-lo com E1a (Tabela~\ref{tab:res-e1a}).
  \item Executar E1b, E2 (Protocolos~S e~U), as ablações e as linhas de base, e preencher a
        Seção~\ref{sec:resultados}.
  \item Executar a Etapa~II com um modelo de linguagem, registrando modelo, versão, tempo e custo.
  \item Revisar a Seção~\ref{sec:critica} e estas conclusões à luz dos resultados.
  \item Preparar a apresentação remota (cronograma proposto no Apêndice~\ref{apx:cronograma}).
\end{enumerate}

\paragraph{Trabalhos futuros.} (i) Combinar o detector com uma regra explícita para atributos que saem de
um valor constante, já que o bit de novidade tem peso apenas $1/N$; (ii) testar uma regra de Hebb centrada para
bits viesados; (iii) usar a rede moderna com memória grande e janelas temporais; (iv) coletar telemetria própria
em baixo nível (por exemplo, com \texttt{psutil} ou eBPF) e injetar falhas controladas (por exemplo, com
\texttt{stress-ng}); (v) reajustar periodicamente o modelo diante da deriva.


\clearpage
\phantomsection
\addcontentsline{toc}{section}{Referências}
\begin{thebibliography}{99}
\setlength{\itemsep}{3pt}

\bibitem{bruck1990} BRUCK, J. On the convergence properties of the Hopfield model.
\textbf{Proceedings of the IEEE}, v.~78, n.~10, p.~1579--1585, 1990.
DOI: \url{https://doi.org/10.1109/5.58341}.

\bibitem{bruckroychowdhury1990} BRUCK, J.; ROYCHOWDHURY, V. P. On the number of spurious memories in the
Hopfield model. \textbf{IEEE Transactions on Information Theory}, v.~36, n.~2, p.~393--397, 1990.
DOI: \url{https://doi.org/10.1109/18.52486}.

\bibitem{chandola2009} CHANDOLA, V.; BANERJEE, A.; KUMAR, V. Anomaly detection: a survey.
\textbf{ACM Computing Surveys}, v.~41, n.~3, art.~15, 2009.
DOI: \url{https://doi.org/10.1145/1541880.1541882}.

\bibitem{demircigil2017} DEMIRCIGIL, M. \textit{et al.} On a model of associative memory with huge storage
capacity. \textbf{Journal of Statistical Physics}, v.~168, n.~2, p.~288--299, 2017.
DOI: \url{https://doi.org/10.1007/s10955-017-1806-y}.

\bibitem{fawcett2006} FAWCETT, T. An introduction to ROC analysis. \textbf{Pattern Recognition Letters},
v.~27, n.~8, p.~861--874, 2006. DOI: \url{https://doi.org/10.1016/j.patrec.2005.10.010}.

\bibitem{hertz1991} HERTZ, J.; KROGH, A.; PALMER, R. G. \textbf{Introduction to the Theory of Neural
Computation}. Redwood City: Addison-Wesley, 1991.

\bibitem{hopfield1982} HOPFIELD, J. J. Neural networks and physical systems with emergent collective
computational abilities. \textbf{Proceedings of the National Academy of Sciences (PNAS)}, v.~79, n.~8,
p.~2554--2558, 1982. DOI: \url{https://doi.org/10.1073/pnas.79.8.2554}.

\bibitem{hopfield1984} HOPFIELD, J. J. Neurons with graded response have collective computational properties
like those of two-state neurons. \textbf{Proceedings of the National Academy of Sciences (PNAS)}, v.~81,
n.~10, p.~3088--3092, 1984. DOI: \url{https://doi.org/10.1073/pnas.81.10.3088}.

\bibitem{hopfieldtank1985} HOPFIELD, J. J.; TANK, D. W. ``Neural'' computation of decisions in optimization
problems. \textbf{Biological Cybernetics}, v.~52, n.~3, p.~141--152, 1985.
DOI: \url{https://doi.org/10.1007/BF00339943}.

\bibitem{kim2022} KIM, S.; CHOI, K.; CHOI, H.-S.; LEE, B.; YOON, S. Towards a rigorous evaluation of
time-series anomaly detection. \textbf{Proceedings of the AAAI Conference on Artificial Intelligence},
v.~36, n.~7, p.~7194--7201, 2022. DOI: \url{https://doi.org/10.1609/aaai.v36i7.20680}.

\bibitem{krotov2016} KROTOV, D.; HOPFIELD, J. J. Dense associative memory for pattern recognition. In:
\textbf{Advances in Neural Information Processing Systems 29 (NeurIPS 2016)}, 2016, p.~1172--1180.
arXiv:1606.01164.

\bibitem{mceliece1987} MCELIECE, R. J.; POSNER, E. C.; RODEMICH, E. R.; VENKATESH, S. S. The capacity of the
Hopfield associative memory. \textbf{IEEE Transactions on Information Theory}, v.~33, n.~4, p.~461--482,
1987. DOI: \url{https://doi.org/10.1109/TIT.1987.1057328}.

\bibitem{hopfieldlayers} ML-JKU. \textbf{hopfield-layers}: código-fonte das camadas Hopfield, indicado por
Ramsauer \textit{et al.} (2021). Disponível em: \url{https://github.com/ml-jku/hopfield-layers}.

\bibitem{omnianomaly} NETMANAIOPS. \textbf{OmniAnomaly}: código e \textit{Server Machine Dataset} (SMD).
Repositório do GitHub (licença MIT declarada no repositório). Disponível em:
\url{https://github.com/NetManAIOps/OmniAnomaly}. Acesso em: 6 nov. 2026.

\bibitem{pedregosa2011} PEDREGOSA, F. \textit{et al.} Scikit-learn: machine learning in Python.
\textbf{Journal of Machine Learning Research}, v.~12, p.~2825--2830, 2011.

\bibitem{ramsauer2021} RAMSAUER, H. \textit{et al.} Hopfield networks is all you need. In:
\textbf{International Conference on Learning Representations (ICLR)}, 2021. arXiv:2008.02217.

\bibitem{saito2015} SAITO, T.; REHMSMEIER, M. The precision-recall plot is more informative than the ROC plot
when evaluating binary classifiers on imbalanced datasets. \textbf{PLOS ONE}, v.~10, n.~3, e0118432, 2015.
DOI: \url{https://doi.org/10.1371/journal.pone.0118432}.

\bibitem{storkey1997} STORKEY, A. J. Increasing the capacity of a Hopfield network without sacrificing
functionality. In: \textbf{Artificial Neural Networks --- ICANN'97}. Lecture Notes in Computer Science.
Berlin: Springer, 1997, p.~451--456. DOI: \url{https://doi.org/10.1007/BFb0020196}.

\bibitem{su2019} SU, Y. \textit{et al.} Robust anomaly detection for multivariate time series through
stochastic recurrent neural network. In: \textbf{Proceedings of the 25th ACM SIGKDD International Conference
on Knowledge Discovery \& Data Mining (KDD '19)}. ACM, 2019.
DOI: \url{https://doi.org/10.1145/3292500.3330672}.

\bibitem{enunciado} TRABALHO de Implementação de Redes Neurais Artificiais: descrição do trabalho. Documento
da disciplina fornecido pelo professor (PDF), [s.d.].

\bibitem{wikihopfield} WIKIPEDIA. \textbf{Hopfield network}. Material fornecido pelo professor (impressão de
14 out. 2026; versão \texttt{oldid=1377477270}). Disponível em:
\url{https://en.wikipedia.org/wiki/Hopfield_network}.

\end{thebibliography}


\appendix
\renewcommand{\thesection}{\Alph{section}}
\section{Pseudocódigos e notas de implementação}\label{apx:algoritmos}

Os algoritmos usam a notação da Seção~\ref{sec:teoria}: $N$ neurônios, $K$ padrões, estados em $\{-1,+1\}^N$ e
limiares $\theta=0$. Descrevem o projeto da implementação; o código dos experimentos não faz parte desta versão
do relatório.

\begin{algorithm}[htbp]
\caption{Armazenamento de $K$ padrões (regra de Hebb ou de Storkey).}\label{alg:armazenar}
\begin{algorithmic}[1]
\Require padrões $\xi^1,\dots,\xi^K\in\{-1,+1\}^N$; regra $\in\{\text{Hebb},\text{Storkey}\}$
\Ensure $W\in\R^{N\times N}$ simétrica, com $w_{ii}=0$
\State $W\gets 0$
\For{$\nu=1,\dots,K$}
  \If{regra $=$ Hebb}
    \State $W\gets W+\tfrac1N\,\xi^\nu\xi^{\nu\top}$ \Comment{Eq.~\eqref{eq:hebb}}
  \Else
    \State $u\gets W\xi^\nu$
    \State $W\gets W+\tfrac1N\bigl(\xi^\nu\xi^{\nu\top}-\xi^\nu u^{\top}-u\,\xi^{\nu\top}+2W\bigr)$
    \Comment{Eq.~\eqref{eq:storkey}}
  \EndIf
  \State zerar a diagonal de $W$
\EndFor
\State \Return $W$
\end{algorithmic}
\end{algorithm}

\paragraph{Forma matricial da regra de Storkey.} Como $w_{ii}=0$, o campo local que exclui $i$ e $j$ é
$h_{ij}^{\nu}=\sum_{k\neq i,j}w_{ik}\xi_k^\nu=u_i-w_{ij}\xi_j^\nu$, com $u=W^{\nu-1}\xi^\nu$, e
analogamente $h_{ji}^{\nu}=u_j-w_{ji}\xi_i^\nu$. Usando $(\xi_i^\nu)^2=(\xi_j^\nu)^2=1$ e $w_{ij}=w_{ji}$,
\[
\xi_i^\nu h_{ji}^\nu+\xi_j^\nu h_{ij}^\nu=\xi_i^\nu u_j+u_i\xi_j^\nu-2w_{ij}^{\nu-1},
\]
e a recursão~\eqref{eq:storkey} fica, para $i\neq j$,
$w_{ij}^{\nu}=w_{ij}^{\nu-1}+\tfrac1N\bigl(\xi_i^\nu\xi_j^\nu-\xi_i^\nu u_j-u_i\xi_j^\nu+2w_{ij}^{\nu-1}\bigr)$,
que é a forma matricial do Algoritmo~\ref{alg:armazenar} fora da diagonal; a diagonal é zerada depois de cada
padrão. Conferimos a equivalência numericamente: com $N=40$ e $K=6$ padrões aleatórios, a diferença máxima entre
a recursão elemento a elemento e a forma matricial foi da ordem de $10^{-16}$.

\begin{algorithm}[htbp]
\caption{Recuperação assíncrona (rede clássica).}\label{alg:recuperar}
\begin{algorithmic}[1]
\Require $W$ (com $\theta=0$); estado inicial $s^0\in\{-1,+1\}^N$; número máximo de varreduras $T_{\max}$
\Ensure estado final $\sstar$; número $T$ de varreduras com ao menos uma troca
\State $s\gets s^0$; \quad $T\gets0$
\For{$t=1,\dots,T_{\max}$}
  \State sortear uma permutação $\pi$ de $\{1,\dots,N\}$; \quad $\mathrm{mudou}\gets\text{falso}$
  \ForAll{$k$ na ordem de $\pi$}
    \State $h\gets\sum_{j}w_{kj}s_j$
    \State $s_k'\gets+1$ se $h\ge0$ e $s_k'\gets-1$ caso contrário \Comment{Eq.~\eqref{eq:update}}
    \If{$s_k'\neq s_k$}
      \State $s_k\gets s_k'$; \quad $\mathrm{mudou}\gets\text{verdadeiro}$
    \EndIf
  \EndFor
  \If{não $\mathrm{mudou}$}
    \State \textbf{interromper}
  \EndIf
  \State $T\gets T+1$
\EndFor
\State \Return $(s,\,T)$
\end{algorithmic}
\end{algorithm}

\begin{algorithm}[htbp]
\caption{Detector \emph{one-class} com a rede de Hopfield clássica.}\label{alg:detectar}
\begin{algorithmic}[1]
\Require $S_{\mathrm{aj}}$ e $S_{\mathrm{cal}}$ (treino sem rótulos); níveis $b$; número $K$ de protótipos;
regra de armazenamento; taxa $\alpha$; $T_{\max}$; escore $a\in\{S_1,\dots,S_6\}$ (Tabela~\ref{tab:escores})
\Ensure para cada observação $x$, o escore $a(x)$ e o rótulo previsto $\hat y$
\Statex \textbf{Ajuste (uma vez por máquina)}
\State determinar os limiares do termômetro e os atributos constantes em $S_{\mathrm{aj}}$
\Comment{Seção~\ref{sec:codificacao}}
\State escolher $K$ medoides $\xi^1,\dots,\xi^K$ entre os códigos de $S_{\mathrm{aj}}$, no espaço de Hamming
\Comment{Seção~\ref{sec:prototipos}}
\State $W\gets$ Algoritmo~\ref{alg:armazenar}$(\xi^1,\dots,\xi^K)$
\State $\tau\gets$ quantil de ordem $1-\alpha$ de $\{a(x):x\in S_{\mathrm{cal}}\}$, com $a(x)$ calculado como abaixo
\Statex \textbf{Decisão (por observação $x$)}
\State $s^0\gets\mathrm{cod}(x)$
\State $(\sstar,T)\gets$ Algoritmo~\ref{alg:recuperar}$(W,s^0,T_{\max})$
\State calcular $a(x)$ a partir de $s^0$, $\sstar$, $T$ e dos protótipos
\State \Return $a(x)$ e $\hat y=\mathbf 1\bigl[a(x)>\tau\bigr]$
\end{algorithmic}
\end{algorithm}

\begin{algorithm}[htbp]
\caption{Recuperação e escores da rede de Hopfield moderna.}\label{alg:moderna}
\begin{algorithmic}[1]
\Require memória $X\in\R^{N\times K}$ (padrões nas colunas); $M=\max_i\|x_i\|$; $\beta>0$; consulta $\xi^0\in\R^N$;
número de atualizações $T\ge1$ ($T=1$ por padrão)
\Ensure estado final $\xi^{T}$; energia $E(\xi^0)$; resíduo $r$; entropia $H$
\State $p\gets\softmax\bigl(\beta X^{\top}\xi^0\bigr)$
\State $E_0\gets-\lse\bigl(\beta,X^{\top}\xi^0\bigr)+\tfrac12\xi^{0\top}\xi^0+\beta^{-1}\ln K+\tfrac12M^2$
\Comment{Eq.~\eqref{eq:energia-moderna}}
\State $H\gets-\sum_ip_i\ln p_i$
\For{$t=1,\dots,T$}
  \State $\xi^{t}\gets X\,\softmax\bigl(\beta X^{\top}\xi^{t-1}\bigr)$ \Comment{Eq.~\eqref{eq:update-moderna}}
\EndFor
\State $r\gets\|\xi^{0}-\xi^{1}\|$
\State \Return $\xi^{T}$, $E_0$, $r$, $H$
\end{algorithmic}
\end{algorithm}

\paragraph{Notas de implementação.}
\begin{itemize}
  \item \textbf{Estabilidade numérica (rede moderna):} $\lse$ e $\softmax$ são calculados subtraindo o máximo de
        $\beta X^{\top}\xi$, para evitar estouro numérico com $\beta$ grande.
  \item \textbf{Vetorização:} a regra de Storkey e as distâncias de Hamming aos protótipos usam produtos
        matriciais; em estados $\pm1$, $\hamm(a,b)=\tfrac12(N-a^{\top}b)$.
  \item \textbf{Ordem de atualização:} a ordem aleatória de cada varredura é gerada por uma semente fixada, e
        os resultados são repetidos com 3 sementes (Seção~\ref{sec:divisao}).
  \item \textbf{Empates:} $h_k=0$ leva a $+1$, conforme a Eq.~\eqref{eq:update}; a Proposição~\ref{prop:conv} cobre
        esse caso.
\end{itemize}

\begin{table}[htbp]
\centering
\caption{Custo aproximado por operação ($T$: número de varreduras ou atualizações; $K$: padrões armazenados).}
\label{tab:complexidade}
\small
\begin{tabular}{@{}lll@{}}
\toprule
Operação & Clássica (Hebb/Storkey) & Moderna\\
\midrule
Armazenar $K$ padrões & $O(KN^2)$ & $O(NK)$ (guardar $X$)\\
Memória & $O(N^2)$ & $O(NK)$\\
Uma varredura / atualização & $O(N^2)$ & $O(NK)$\\
Recuperação completa & $O(TN^2)$ & $O(TNK)$\\
Escores de distância aos protótipos & $O(NK)$ & ---\\
\bottomrule
\end{tabular}
\end{table}

\section{Exemplo numérico com 5 neurônios}\label{apx:exemplo}

O exemplo abaixo, calculado à mão e conferido numericamente, mostra a recuperação e os escores da
Tabela~\ref{tab:escores} em uma rede com um único padrão armazenado, $\xi=(+1,-1,+1,-1,+1)$ e $N=5$. Para
simplificar, $W=\xi\xi^{\top}-I$ (regra de Hebb sem o fator $1/N$ da Eq.~\eqref{eq:hebb}, o que não altera a
dinâmica com $\theta=0$; com o fator, todas as energias ficam divididas por $N=5$). Assim
$w_{ij}=\xi_i\xi_j$ para $i\neq j$ e $w_{ii}=0$.

\paragraph{Energia em função da sobreposição.} Como $\sum_{i,j}\xi_i\xi_js_is_j=(\xi\cdot s)^2$ e os termos
diagonais somam $N$,
\[
E(s)=-\tfrac12 s^{\top}Ws=-\tfrac12\bigl[(\xi\cdot s)^2-N\bigr].
\]
O padrão armazenado tem $\xi\cdot\xi=5$ e $E(\xi)=-\tfrac12(25-5)=-10$.

\paragraph{Consulta com um bit diferente.} Seja $s^0=(+1,-1,-1,-1,+1)$, que difere de $\xi$ no neurônio~3, com
$\xi\cdot s^0=3$ e $E(s^0)=-\tfrac12(9-5)=-2$. O campo local de cada neurônio é
$h_i=\sum_{j\neq i}\xi_i\xi_js_j=\xi_i(\xi\cdot s)-s_i$:
\begin{center}
\small
\begin{tabular}{@{}lccccc@{}}
\toprule
Neurônio $i$ & 1 & 2 & 3 & 4 & 5\\
\midrule
$\xi_i$ & $+1$ & $-1$ & $+1$ & $-1$ & $+1$\\
$s^0_i$ & $+1$ & $-1$ & $-1$ & $-1$ & $+1$\\
$h_i=3\xi_i-s^0_i$ & $2$ & $-2$ & $4$ & $-2$ & $2$\\
$\sgn(h_i)$ & $+1$ & $-1$ & $+1$ & $-1$ & $+1$\\
$s^{\ast}_i$ & $+1$ & $-1$ & $+1$ & $-1$ & $+1$\\
\bottomrule
\end{tabular}
\end{center}
Só o neurônio~3 muda: $s^0_3=-1$ e $h_3=4>0$. Qualquer que seja a ordem de visita, o resultado é o mesmo:
$\sstar=\xi$, em uma varredura com troca (mais uma sem troca, que confirma o ponto fixo). A variação de energia
da troca é $-(s_3'-s_3)h_3=-2\cdot4=-8$, de $-2$ para $-10$, como previsto pela demonstração da
Proposição~\ref{prop:conv}.

\paragraph{Consulta igual ao padrão invertido.} Com $s^0=-\xi$, $\xi\cdot s^0=-5$ e
$h_i=-5\xi_i-s^0_i=-4\xi_i=4s^0_i$: cada neurônio já concorda com o sinal do seu campo, e $-\xi$ é um ponto fixo
com $E(-\xi)=-10=E(\xi)$, porque a energia é par em $s$ quando $\theta=0$. A rede não reconhece nada de errado;
só os escores que medem a distância \emph{com sinal} aos protótipos, $S_2$ e $S_5$, denunciam o padrão invertido.

\begin{table}[htbp]
\centering
\caption{Escores do exemplo (Tabela~\ref{tab:escores}). Quanto maior, mais anômalo; $S_3$ e $S_4$ têm a escala da
energia sem o fator $1/N$.}
\label{tab:ex-escores}
\small
\begin{tabular}{@{}lrrrrrr@{}}
\toprule
Consulta $s^0$ & $S_1$ & $S_2$ & $S_3$ & $S_4$ & $S_5$ & $S_6$\\
\midrule
$\xi$ (padrão armazenado) & 0 & 0 & $-10$ & 0 & 0 & 0\\
$(+1,-1,-1,-1,+1)$ (1 bit diferente) & 0,2 & 0,2 & $-2$ & 8 & 0 & 1\\
$-\xi$ (padrão invertido) & 0 & 1 & $-10$ & 0 & 1 & 0\\
\bottomrule
\end{tabular}
\end{table}

\noindent A Tabela~\ref{tab:ex-escores} resume a lição: $S_1$, $S_3$ e $S_4$ (que dependem só da dinâmica e da energia) dão
valores indistinguíveis dos do padrão armazenado para $-\xi$; por isso o detector usa também $S_2$ e $S_5$.

\section{\emph{Prompt} e registro da Etapa II}\label{apx:prompt}

O enunciado pede que o \emph{prompt} e a versão do modelo sejam documentados~\cite{enunciado}. Esta seção fixa o
\emph{prompt} antes da execução. O texto é em inglês e usa apenas caracteres ASCII, para evitar problemas de
codificação em APIs; os campos entre chaves são preenchidos pelo código.

\subsection{Mensagens}
\begin{lstlisting}
=== SYSTEM MESSAGE ===
You are an expert in server monitoring and anomaly detection.

You will receive one observation of a server: 38 anonymous
resource-utilization and activity metrics, named m1 to m38. Every value is
scaled to the range [0, 1]. Metrics are sampled once per minute. Their names
and units are not available.

Task: decide whether the observation belongs to an anomalous period (a server
incident) or to normal operation.

Rules:
- Use only the information given in this conversation.
- Anomalies are rare: roughly 4 out of every 100 observations.
- Reply with one JSON object and nothing else:
  {"label": 0 or 1, "score": a number between 0 and 1}
  label = 1 if the observation is anomalous, 0 if it is normal.
  score = your confidence that the observation is anomalous
          (0 = certainly normal, 1 = certainly anomalous).

=== USER MESSAGE, zero-shot ===
Observation:
{observation}

=== USER MESSAGE, few-shot (k = 3, 5, 10 or 20 labeled examples) ===
Labeled examples from the same kind of server (k = {k}):

Example 1:
{observation_1}
Answer: {"label": 0, "score": 0.0}

Example 2:
{observation_2}
Answer: {"label": 1, "score": 1.0}

... (up to Example {k}, in random order)

Now classify this observation:
{observation}

=== OPTIONAL BLOCK, reference profile (zero-shot variant) ===
Typical values of each metric during normal operation of this server
(median / 95th percentile):
m1: {median_1} / {p95_1}
...
m38: {median_38} / {p95_38}
\end{lstlisting}

\subsection{Serialização das observações}
Cada observação é escrita em uma linha, com os 38 valores no formato \texttt{mj=v}, separados por espaço e com
4 casas decimais. O exemplo abaixo é a linha 15.000 do arquivo de teste da máquina 1-1 (rótulo 0), quebrada em cinco linhas apenas para caber na página; no \emph{prompt}, é uma linha só. A linha 16.100 da mesma máquina
(rótulo 1) começa com \texttt{m1=0.5161}, acima do máximo de treino do atributo (0,4946), e serve de exemplo de
observação anômala (Seção~\ref{sec:exemplos}). Ambas são ilustrações da \emph{forma} do texto, não exemplos do
protocolo.
\begin{lstlisting}
m1=0.0968 m2=0.0794 m3=0.0702 m4=0.0627 m5=0.0000 m6=0.8974 m7=0.2750 m8=0.0000
m9=0.0288 m10=0.0000 m11=0.1066 m12=0.1081 m13=0.0274 m14=0.0575 m15=0.1250 m16=0.1246
m17=0.0000 m18=0.0000 m19=0.1483 m20=0.1009 m21=0.1254 m22=0.0984 m23=0.8330 m24=0.0126
m25=0.0493 m26=0.0117 m27=0.0000 m28=0.1029 m29=0.0000 m30=0.0201 m31=0.0897 m32=0.0839
m33=0.0004 m34=0.0001 m35=0.0810 m36=0.0811 m37=0.0000 m38=0.0000
\end{lstlisting}

\subsection{Regras de execução}
\begin{itemize}
  \item \textbf{Exemplos:} para cada par $(k,\text{semente})$ sorteia-se de $L_{\mathrm{tr}}$ uma lista fixa de $k$
        exemplos, com $k=3$: 2~normais e 1~anômalo; $k=5$: 3+2; $k=10$: 5+5; $k=20$: 10+10, em ordem
        aleatória. A mesma lista é usada para todas as observações da execução, o que permite o uso de
        \emph{prompt caching}, se a API o oferecer.
  \item \textbf{Parâmetros:} temperatura~0 (ou a menor disponível) e limite curto de \emph{tokens} de saída; usar o
        modo de saída JSON quando existir.
  \item \textbf{Respostas inválidas:} se a resposta não for um JSON válido, repete-se a chamada uma vez; persistindo
        o erro, a observação conta como rótulo~0 e escore~0, e a taxa de respostas inválidas é reportada.
  \item \textbf{Escore:} o escore autodeclarado não é calibrado, de modo que a AUC-ROC calculada sobre ele é
        apenas indicativa. Se a API fornecer \emph{logprobs}, a probabilidade do \emph{token} do rótulo é uma
        alternativa.
  \item \textbf{Número de chamadas:} para $n$ observações por máquina, são $3\times[1+4\times3]\times n=39\,n$
        chamadas (3 máquinas, 1 configuração \emph{zero-shot} e 4 valores de $k$ com 3 sementes cada); por
        exemplo, 19.500 chamadas com $n=500$. O custo total segue de multiplicar esse número pelos \emph{tokens}
        médios por chamada e pelos preços da Tabela~\ref{tab:registro-llm}.
  \item \textbf{Contaminação:} o SMD é público desde 2019~\cite{su2019,omnianomaly}, e não é possível excluir
        que faça parte dos dados de pré-treinamento do modelo; isso deve ser mencionado na discussão.
\end{itemize}

\subsection{Registro da execução}
\begin{table}[htbp]
\centering
\caption{Registro da Etapa II (preencher durante a execução).}
\label{tab:registro-llm}
\small
\begin{tabularx}{\textwidth}{@{}lL@{}}
\toprule
Campo & Valor\\
\midrule
Provedor e modelo & \apreencher\\
Identificador exato da versão (cadeia da API) & \apreencher\\
Data(s) de execução & \apreencher\\
Temperatura e limite de \emph{tokens} de saída & 0 e valor curto (proposta); valores usados: \apreencher\\
Observações avaliadas por máquina ($n$) & 500 a 2.000 (proposta); valor usado: \apreencher\\
\emph{Tokens} médios por chamada (entrada / saída) & \apreencher\\
Preço por milhão de \emph{tokens} (entrada / saída) & \apreencher\\
Tempo por 100 amostras & \apreencher\\
Custo estimado por 100 amostras & \apreencher\\
Respostas inválidas (\%) & \apreencher\\
\bottomrule
\end{tabularx}
\end{table}

\section{Cronograma proposto}\label{apx:cronograma}

As apresentações remotas ocorrem em 02 e 04/12~\cite{enunciado}. O enunciado não informa a data de entrega do
relatório, nem em qual dos dois dias cada equipe apresenta; o plano abaixo é uma \textbf{proposta}, que supõe a
entrega do relatório e dos \emph{slides} até a véspera da primeira apresentação (01/12), e deve ser confirmada com
o professor.

\begin{table}[htbp]
\centering
\caption{Cronograma proposto (datas de 2026).}
\label{tab:cronograma}
\small
\begin{tabularx}{\textwidth}{@{}c>{\raggedright\arraybackslash}p{2.3cm}>{\hsize=1.5\hsize}L>{\hsize=.5\hsize}L@{}}
\toprule
 & Período & Atividades & Entrega\\
\midrule
1 & 06--13/11 &
Confirmar com o professor o uso da rede de Hopfield como modelo do tema, as linhas de base e as datas;
dividir as tarefas; baixar o SMD (registrar o \emph{commit}); implementar leitura, divisões e codificação em
termômetro; implementar Hebb, Storkey, recuperação e escores; executar E1a. &
Dados e rede clássica validados em E1a\\
2 & 16--20/11 (20/11: feriado) &
Seleção de protótipos e E1b; detector nas máquinas 1-3, 2-5 e 3-6; linhas de base (\emph{Random Forest}, SVM,
OC-SVM e \emph{k}-NN); escolha do limiar em $L_{\mathrm{val}}$. &
Primeiros resultados da Etapa~I nas 3 máquinas\\
3 & 23--27/11 &
Protocolo~S (18 máquinas) e Protocolo~U (28); ablações; rede moderna ($\beta$); Etapa~II (\emph{zero-shot} e
\emph{few-shot}, com registro de tempo e custo); preencher a Seção~\ref{sec:resultados}. &
Seção~\ref{sec:resultados} completa\\
4 & 30/11--01/12 &
Revisar as Seções~\ref{sec:critica} e~\ref{sec:conclusoes} e o resumo à luz dos resultados; conferir as
referências; preparar os \emph{slides} e ensaiar. &
Relatório e \emph{slides}\\
5 & 02 e 04/12 & Apresentação remota. & ---\\
\bottomrule
\end{tabularx}
\end{table}

\paragraph{Divisão de tarefas (equipe de até 3 integrantes).}
\begin{itemize}
  \item \textbf{Integrante A} (\apreencher[nome]): dados, divisões, linhas de base clássicas e métricas.
  \item \textbf{Integrante B} (\apreencher[nome]): rede de Hopfield (clássica e moderna), E1 e ablações.
  \item \textbf{Integrante C} (\apreencher[nome]): Etapa~II (\emph{prompt}, execução e custo) e revisão do texto.
\end{itemize}

\paragraph{Riscos e mitigação.}
\begin{itemize}
  \item \textbf{Hopfield próximo do acaso:} o resultado é registrado e analisado como achado (ablações e
        Seção~\ref{sec:critica}); a comparação crítica vale igualmente.
  \item \textbf{Prazo apertado:} priorizar o Protocolo~S nas 3 máquinas representativas e estender às 18 depois.
  \item \textbf{Custo ou latência da API:} reduzir $n$, usar o modelo mais barato disponível e limitar $k=20$ a
        uma semente.
  \item \textbf{Modelo fora da lista do enunciado:} a confirmação com o professor é a primeira atividade do
        cronograma.
\end{itemize}

\section{Conformidade com o enunciado}\label{apx:checklist}

\newcommand{\stpronto}{\textcolor{green!45!black}{\textbf{Pronto}}}
\newcommand{\stparcial}{\textcolor{orange!80!black}{\textbf{Parcial}}}
\newcommand{\stpendente}{\textcolor{red!70!black}{\textbf{Pendente}}}

A Tabela~\ref{tab:checklist} relaciona cada exigência do enunciado~\cite{enunciado} ao local do relatório e ao seu
estado. \stpronto: texto ou dados concluídos nesta versão; \stparcial: método, critérios e formato de relato
fixados, mas a evidência depende dos experimentos; \stpendente: ainda a fazer.

{\small
\begin{longtable}{@{}p{5.4cm}p{3.7cm}p{5.4cm}@{}}
\caption{Exigências do enunciado, local no relatório e estado.}\label{tab:checklist}\\
\toprule
Exigência & Onde & Estado\\
\midrule
\endfirsthead
\toprule
Exigência & Onde & Estado\\
\midrule
\endhead
\midrule
\multicolumn{3}{r}{\emph{continua na próxima página}}\\
\endfoot
\bottomrule
\endlastfoot
Escolha do problema e da base de dados; tipo (classificação) &
Seções~\ref{sec:problema} e~\ref{sec:smd} &
\stpronto: classificação binária por ponto no SMD\\[3pt]
Definições do problema e do conjunto de treinamento &
Seções~\ref{sec:problema} e~\ref{sec:smd} &
\stpronto\\[3pt]
2 a 3 exemplos da base de dados &
Seção~\ref{sec:exemplos} &
\stpronto: 3 observações reais da máquina 1-1\\[3pt]
Valores desejados e atributos &
Seções~\ref{sec:problema} e~\ref{sec:smd} &
\stpronto: rótulo 0/1 e 38 atributos\\[3pt]
Quantidade de amostras &
Seção~\ref{sec:smd}; Apêndice~\ref{apx:maquinas} &
\stpronto\\[3pt]
Preparação e limpeza &
Seção~\ref{sec:prep} &
\stpronto\\[3pt]
Balanceamento de classes &
Seção~\ref{sec:balanceamento} &
\stpronto\\[3pt]
Divisão dos conjuntos &
Seção~\ref{sec:divisao} &
\stpronto\\[3pt]
Métricas de avaliação &
Seção~\ref{sec:metricas} &
\stpronto\\[3pt]
Etapa~I: ao menos 2 modelos de classificação &
Seção~\ref{sec:baselines}; Seção~\ref{sec:resultados} &
\stparcial: protocolo definido, execução pendente; o Hopfield não consta da lista do enunciado (confirmar com o
professor)\\[3pt]
Métricas mínimas: acurácia, precisão, revocação, $F$-scores, AUC-ROC e matriz de confusão &
Seção~\ref{sec:metricas}; Tabelas~\ref{tab:res-s} e~\ref{tab:res-conf} &
\stparcial: formato fixado, valores pendentes\\[3pt]
Etapa~II: classificação por \emph{prompting} (descrição da tarefa, \emph{zero-shot}, \emph{few-shot} com
$\{3,5,10,20\}$) &
Seção~\ref{sec:etapa2}; Apêndice~\ref{apx:prompt} &
\stparcial: \emph{prompt} documentado, execução pendente\\[3pt]
Etapa~II: repetir as métricas; tempo por 100 amostras e custo estimado &
Seção~\ref{sec:etapa2}; Tabelas~\ref{tab:res-llm} e~\ref{tab:registro-llm} &
\stparcial: formato fixado, valores pendentes\\[3pt]
Mesmos dados da Etapa~I; documentar \emph{prompt} e versão do modelo &
Seção~\ref{sec:etapa2}; Apêndice~\ref{apx:prompt} &
\stparcial: \emph{prompt} documentado, versão do modelo pendente\\[3pt]
Comparação crítica (pontos positivos e negativos) &
Seção~\ref{sec:critica} &
\stparcial: critérios e análise \emph{a priori} prontos, evidência empírica pendente\\[3pt]
Conclusões: o que funciona, limites e próximos passos &
Seção~\ref{sec:conclusoes} &
\stparcial: limites e próximos passos prontos; ``o que funciona'' depende dos experimentos\\[3pt]
Equipe de no máximo 3 integrantes &
Folha de rosto; Apêndice~\ref{apx:cronograma} &
\stpendente: nomes e instituição\\[3pt]
Relatório técnico &
Este documento &
\stparcial: estrutura completa, resultados pendentes\\[3pt]
Apresentação remota (02 e 04/12) &
Apêndice~\ref{apx:cronograma} &
\stpendente\\
\end{longtable}
}

\section{Derivações}\label{apx:deriv}

As derivações abaixo são nossas (a menos de onde há citação) e complementam as seções~\ref{sec:teoria}
e~\ref{sec:metodo}.

\subsection{Energia de Hebb em função das sobreposições}
Com $w_{ij}=\tfrac1N\sum_{\mu}\xi_i^\mu\xi_j^\mu$ para $i\neq j$, $w_{ii}=0$, $\theta=0$ e a sobreposição
$m_\mu(s)=\tfrac1N\,\xi^\mu\!\cdot s$, como $(\xi_i^\mu s_i)^2=1$,
\[
E(s)=-\tfrac12\sum_{i\neq j}w_{ij}s_is_j
=-\frac1{2N}\sum_{\mu=1}^{K}\Bigl[\Bigl(\sum_i\xi_i^\mu s_i\Bigr)^{2}-\sum_i(\xi_i^\mu s_i)^2\Bigr]
=-\frac N2\sum_{\mu=1}^{K}m_\mu(s)^2+\frac K2 .
\]
Consequências: (i) $E$ depende de $s$ só pelas sobreposições, e seus mínimos são estados com $|m_\mu|$ grande;
(ii) $E(-s)=E(s)$, de modo que $-\xi^\mu$ é tão profundo quanto $\xi^\mu$; (iii) se os outros padrões são
quase ortogonais a $\xi^\nu$, $E(\xi^\nu)\approx-\tfrac N2+\tfrac K2$. A energia da CAM de Hopfield e
Tank~\cite[Eq.~5]{hopfieldtank1985}, $-\tfrac12\sum_s(V^s\!\cdot V)^2$, é $N$ vezes a nossa a menos de uma
constante aditiva, com mínimos de profundidade ${\approx}-\tfrac12N^2$.

\subsection{Ruído cruzado da regra de Hebb}
Para o padrão armazenado $\xi^\nu$, o campo local se decompõe em
$h_i(\xi^\nu)=\tfrac{N-1}{N}\xi_i^\nu+\eta_i$, com
$\eta_i=\tfrac1N\sum_{\mu\neq\nu}\sum_{j\neq i}\xi_i^\mu\xi_j^\mu\xi_j^\nu$. Se os padrões são independentes e não
viesados, cada parcela $\xi_i^\mu\xi_j^\mu\xi_j^\nu$ ($\mu\neq\nu$, $j\neq i$) é uma variável $\pm1$ uniforme e
independente das demais (contém o fator $\xi_j^\mu$, exclusivo do par $(\mu,j)$). Logo $\mathbb E[\eta_i]=0$ e
$\mathrm{Var}[\eta_i]=\tfrac{(K-1)(N-1)}{N^2}\approx\tfrac{K-1}{N}$. Pelo teorema central do limite, o bit $i$ é
trocado no primeiro passo com probabilidade ${\approx}\,\Phi\bigl(-\sqrt{N/K}\bigr)$, em que $\Phi$ é a função de
distribuição normal padrão. A Tabela~\ref{tab:ruido} mostra a aproximação e o resultado de uma simulação
($N=1000$, 20 repetições). É apenas uma estimativa do \emph{primeiro passo} a partir de um padrão exato; o
limite crítico de $0{,}138N$ vem de uma análise mais completa~\cite{hertz1991}.
\begin{table}[htbp]
\centering
\caption{Probabilidade de erro por bit no primeiro passo, a partir de um padrão armazenado (regra de Hebb,
padrões aleatórios não viesados).}
\label{tab:ruido}
\small
\begin{tabular}{@{}rrrr@{}}
\toprule
$K/N$ & Desvio padrão do ruído $\sqrt{K/N}$ & $\Phi(-\sqrt{N/K})$ & Simulação\\
\midrule
0,05  & 0,22 & $3{,}9\times10^{-6}$ & $5{,}0\times10^{-6}$\\
0,10  & 0,32 & $7{,}8\times10^{-4}$ & $7{,}0\times10^{-4}$\\
0,138 & 0,37 & $3{,}6\times10^{-3}$ & $3{,}5\times10^{-3}$\\
0,20  & 0,45 & $1{,}3\times10^{-2}$ & $1{,}3\times10^{-2}$\\
\bottomrule
\end{tabular}
\end{table}

\subsection{A energia da rede moderna como mínimo suave}
Como $x_i^{\top}\xi=\tfrac12\|x_i\|^2+\tfrac12\|\xi\|^2-\tfrac12\|\xi-x_i\|^2$,
\[
\lse(\beta,X^{\top}\xi)=\beta^{-1}\ln\sum_i e^{\beta x_i^{\top}\xi}
=\tfrac12\|\xi\|^2+\beta^{-1}\ln\sum_i e^{\frac12\beta\|x_i\|^2}e^{-\frac12\beta\|\xi-x_i\|^2}.
\]
Em~\eqref{eq:energia-moderna} o termo $\tfrac12\xi^{\top}\xi$ se cancela, e
\[
E(\xi)=-\beta^{-1}\ln\sum_i e^{\frac12\beta\|x_i\|^2}e^{-\frac12\beta\|\xi-x_i\|^2}+\beta^{-1}\ln K+\tfrac12M^2 .
\]
Se $\|x_i\|=M$ para todo $i$, o fator $e^{\frac12\beta M^2}$ sai da soma e cancela $\tfrac12M^2$. Com
$d_i=\tfrac12\|\xi-x_i\|^2$,
\[
E(\xi)=-\beta^{-1}\ln\Bigl[\frac1K\sum_{i=1}^{K}e^{-\beta d_i}\Bigr],\qquad
\min_i d_i\;\le\;E(\xi)\;\le\;\frac1K\sum_i d_i .
\]
A primeira desigualdade vem de $e^{-\beta d_i}\le e^{-\beta\min_j d_j}$ para todo $i$; a segunda, da
desigualdade de Jensen. Quando $\beta\to\infty$, $E\to\min_i d_i$ (vizinho mais próximo); quando $\beta\to0$,
$E\to\tfrac1K\sum_i d_i$, cujo mínimo em $\xi$ é o centroide dos padrões. Além disso,
$e^{-\beta E(\xi)}=\tfrac1K\sum_i\exp\bigl(-\|\xi-x_i\|^2/(2\sigma^2)\bigr)$ com $\sigma^2=1/\beta$ é, a menos
de uma constante, uma estimativa de densidade por kernel gaussiano, de modo que $E(\xi)=-\beta^{-1}\ln\hat p(\xi)+\text{const.}$

\subsection{Padrões binários: distância de Hamming e separação}
Para $x,y\in\{-1,+1\}^N$, coordenadas iguais contribuem com $+1$ a $x^{\top}y$ e coordenadas diferentes, com
$-1$; logo $x^{\top}y=N-2\hamm(x,y)$, $\|x-y\|^2=4\hamm(x,y)$ e $\tfrac12\|x-y\|^2=2\hamm(x,y)$. Para a separação
$\Delta_i=x_i^{\top}x_i-\max_{j\neq i}x_i^{\top}x_j$ dos Teoremas~4 e~5 de Ramsauer et al.~\cite{ramsauer2021},
\[
\Delta_i=N-\bigl(N-2\min_{j\neq i}\hamm(x_i,x_j)\bigr)=2\min_{j\neq i}\hamm(x_i,x_j).
\]
Portanto, com protótipos binários, a energia moderna é um mínimo suave das distâncias de Hamming,
$E(s)=-\beta^{-1}\ln\bigl[\tfrac1K\sum_ie^{-2\beta\hamm(s,x_i)}\bigr]$, e a facilidade de recuperar um protótipo
depende da sua distância de Hamming ao protótipo mais próximo.

\subsection{Distância de Hamming na codificação em termômetro}
Para um atributo com limiares $\tau_1<\dots<\tau_{b'}$ (distintos, como na Seção~\ref{sec:codificacao}), seja
$r(x)=\#\{l:x>\tau_l\}$ o posto de $x$. Como os limiares são crescentes, $x>\tau_l$ se, e somente se,
$l\le r(x)$, e o código é $(+1)^{r(x)}(-1)^{b'-r(x)}$. Para dois valores com postos $r\le r'$, os códigos diferem
exatamente nas posições $r<l\le r'$, de modo que
\[
\hamm\bigl(\mathrm{cod}(x),\mathrm{cod}(x')\bigr)=\bigl|r(x)-r(x')\bigr| .
\]
A distância entre estados completos é a soma dessas diferenças de postos sobre os atributos, mais as
discordâncias nos bits de novidade. Como os limiares são quantis empíricos do treino, $r(x)/(b+1)$ aproxima a
função de distribuição empírica do atributo; assim a distância é invariante a reescalonamentos monótonos de
cada atributo.

\section{Dados por máquina e divisão cronológica}\label{apx:maquinas}

A Tabela~\ref{tab:maquinas} traz, para cada uma das 28 máquinas do SMD, o tamanho de cada metade, as anomalias
rotuladas, os cortes cronológicos da Seção~\ref{sec:divisao} e a elegibilidade para os protocolos. Todos os
números são de cálculo próprio a partir dos arquivos do repositório~\cite{omnianomaly}.

{\footnotesize
\setlength{\tabcolsep}{3pt}
\begin{longtable}{@{}lrrrrrrrrrrc@{}}
\caption{Dados por máquina e divisão cronológica. \textbf{Linhas:} pontos de 1 minuto de cada metade.
\textbf{Teste:} \emph{Anôm.} é o número de pontos rotulados como anômalos, \emph{\%} a sua proporção e
\emph{Seg.} o número de segmentos anômalos contíguos. \textbf{Cortes:} $c_1=\mathrm{round}(0{,}4n)$ e
$c_2=\mathrm{round}(0{,}6n)$, com $n$ igual ao número de linhas de teste e índices a partir de 0;
$^{\dagger}$: corte deslocado para uma zona sem anomalias. \textbf{Anomalias por bloco:} em $L_{\mathrm{tr}}$,
$L_{\mathrm{val}}$ e $L_{\mathrm{te}}$ (Eq.~\ref{eq:blocos}); a soma dos três blocos é igual a \emph{Anôm.},
isto é, nenhuma anomalia cai nos embargos. \textbf{Elig.:} S = supervisionados e \emph{one-class};
O = só \emph{one-class}; --- = nenhum (Tabela~\ref{tab:elegib}).}
\label{tab:maquinas}\\
\toprule
 & \multicolumn{2}{c}{Linhas} & \multicolumn{3}{c}{Teste} & \multicolumn{2}{c}{Cortes} &
 \multicolumn{3}{c}{Anomalias por bloco} & \\
\cmidrule(lr){2-3}\cmidrule(lr){4-6}\cmidrule(lr){7-8}\cmidrule(lr){9-11}
Máq. & Treino & Teste & Anôm. & \% & Seg. & $c_1$ & $c_2$ & $L_{\mathrm{tr}}$ & $L_{\mathrm{val}}$ &
$L_{\mathrm{te}}$ & Elig.\\
\midrule
\endfirsthead
\multicolumn{12}{l}{\emph{Tabela~\ref{tab:maquinas} (continuação)}}\\
\toprule
 & \multicolumn{2}{c}{Linhas} & \multicolumn{3}{c}{Teste} & \multicolumn{2}{c}{Cortes} &
 \multicolumn{3}{c}{Anomalias por bloco} & \\
\cmidrule(lr){2-3}\cmidrule(lr){4-6}\cmidrule(lr){7-8}\cmidrule(lr){9-11}
Máq. & Treino & Teste & Anôm. & \% & Seg. & $c_1$ & $c_2$ & $L_{\mathrm{tr}}$ & $L_{\mathrm{val}}$ &
$L_{\mathrm{te}}$ & Elig.\\
\midrule
\endhead
\midrule
\multicolumn{12}{r}{\emph{continua na próxima página}}\\
\endfoot
\bottomrule
\endlastfoot
1-1 & 28.479 & 28.479 & 2.694 & 9,46 & 8 & 11.392 & 16.903$^{\dagger}$ & 0 & 546 & 2.148 & O \\
1-2 & 23.694 & 23.694 & 542 & 2,29 & 10 & 9.478 & 14.216 & 140 & 0 & 402 & --- \\
1-3 & 23.702 & 23.703 & 817 & 3,45 & 12 & 9.481 & 14.222 & 86 & 281 & 450 & S \\
1-4 & 23.706 & 23.707 & 720 & 3,04 & 12 & 9.483 & 14.224 & 85 & 279 & 356 & S \\
1-5 & 23.705 & 23.706 & 100 & 0,42 & 7 & 9.482 & 14.224 & 0 & 73 & 27 & O \\
1-6 & 23.688 & 23.689 & 3.708 & 15,65 & 30 & 9.476 & 14.213 & 472 & 27 & 3.209 & S \\
1-7 & 23.697 & 23.697 & 2.398 & 10,12 & 13 & 9.479 & 14.218 & 2.247 & 41 & 110 & S \\
1-8 & 23.698 & 23.699 & 763 & 3,22 & 20 & 9.480 & 14.219 & 477 & 16 & 270 & S \\
\midrule
2-1 & 23.693 & 23.694 & 1.170 & 4,94 & 13 & 9.478 & 14.216 & 122 & 34 & 1.014 & S \\
2-2 & 23.699 & 23.700 & 2.833 & 11,95 & 11 & 9.480 & 14.220 & 0 & 11 & 2.822 & O \\
2-3 & 23.688 & 23.689 & 269 & 1,14 & 10 & 9.476 & 14.213 & 104 & 30 & 135 & S \\
2-4 & 23.689 & 23.689 & 1.694 & 7,15 & 20 & 9.413$^{\dagger}$ & 14.213 & 310 & 67 & 1.317 & S \\
2-5 & 23.688 & 23.689 & 980 & 4,14 & 21 & 9.476 & 14.213 & 314 & 39 & 627 & S \\
2-6 & 28.743 & 28.743 & 424 & 1,48 & 8 & 11.497 & 17.246 & 78 & 128 & 218 & S \\
2-7 & 23.696 & 23.696 & 417 & 1,76 & 20 & 9.478 & 14.218 & 27 & 33 & 357 & S \\
2-8 & 23.702 & 23.703 & 161 & 0,68 & 1 & 9.481 & 14.222 & 0 & 0 & 161 & --- \\
2-9 & 28.722 & 28.722 & 1.755 & 6,11 & 10 & 11.489 & 17.346$^{\dagger}$ & 0 & 992 & 763 & O \\
\midrule
3-1 & 28.700 & 28.700 & 308 & 1,07 & 4 & 11.480 & 17.220 & 183 & 0 & 125 & --- \\
3-2 & 23.702 & 23.703 & 1.109 & 4,68 & 10 & 9.481 & 14.222 & 873 & 127 & 109 & S \\
3-3 & 23.703 & 23.703 & 632 & 2,67 & 26 & 9.481 & 14.222 & 52 & 15 & 565 & S \\
3-4 & 23.687 & 23.687 & 977 & 4,12 & 8 & 9.475 & 14.212 & 865 & 20 & 92 & S \\
3-5 & 23.690 & 23.691 & 426 & 1,80 & 11 & 9.476 & 14.215 & 155 & 3 & 268 & S \\
3-6 & 28.726 & 28.726 & 1.194 & 4,16 & 11 & 11.490 & 17.236 & 174 & 437 & 583 & S \\
3-7 & 28.705 & 28.705 & 434 & 1,51 & 5 & 11.431$^{\dagger}$ & 17.223 & 35 & 81 & 318 & S \\
3-8 & 28.703 & 28.704 & 1.371 & 4,78 & 6 & 11.482 & 17.336$^{\dagger}$ & 0 & 908 & 463 & O \\
3-9 & 28.713 & 28.713 & 303 & 1,06 & 4 & 11.485 & 17.228 & 202 & 0 & 101 & --- \\
3-10 & 23.692 & 23.693 & 1.047 & 4,42 & 13 & 9.477 & 14.216 & 188 & 312 & 547 & S \\
3-11 & 28.695 & 28.696 & 198 & 0,69 & 3 & 11.478 & 17.218 & 0 & 0 & 198 & --- \\
\midrule
Total & 708.405 & 708.420 & 29.444 & 4,16 & 327 & --- & --- & 7.189 & 4.500 & 17.755 & --- \\
\end{longtable}
}


\end{document}
